\subsection{分配性}

定义5. 设 \(\mathbf{E}_1, \ldots, \mathbf{E}_n\) 和 \(\mathbf{F}\) 为集合， \(u\) 为 \(\mathbf{E}_1 \times \cdots \times \mathbf{E}_n\) 到 \(\mathbf{F}\) 中的一个映射。设 \(i \in [1, n]\) 。假设 \(\mathbf{E}_i\) 和 \(\mathbf{F}\) 被赋予广群结构。称 \(u\) 关于指标变量 \(i\) 是分配的，如果部分映射

\[
x _ {i} \mapsto u (a _ {1}, \dots , a _ {i - 1}, x _ {i}, a _ {i + 1}, \dots , a _ {n})
\]

是 \(\mathbf{E}_i\) 到 \(\mathbf{F}\) 的同态，其中 \(a_j\) 是 \(\mathbf{E}_j\) 中任意固定的元素且 \(j \neq i\) 。

若 \(\top\) 表示 \(\mathbf{E}_i\) 和 \(\mathbf{F}\) 上的内合成律，则 \(u\) 的分配性由下列等式给出：

\[
\begin{array}{r l} & u (a _ {1}, \dots , a _ {i - 1}, x _ {i} \top x _ {i} ^ {\prime}, a _ {i + 1}, a _ {n}) \\ & \quad = u (a _ {1}, \dots , a _ {i - 1}, x _ {i}, a _ {i + 1}, \dots , a _ {n}) \top u (a _ {1}, \dots , a _ {i - 1}, x _ {i} ^ {\prime}, a _ {i + 1}, \dots , a _ {n}) \end{array}\tag{1}
\]

对于 \(i = 1, 2, \ldots, n\) , \(a_1 \in \mathrm{E}_1, \ldots, a_{i-1} \in \mathrm{E}_{i-1}\) , \(x_i \in \mathrm{E}_i\) , \(x_i' \in \mathrm{E}_i\) , \(a_{i+1} \in \mathrm{E}_{i+1}, \ldots, a_n \in \mathrm{E}_n\) .

例。设 E 是乘法记法书写的幺半群（相应地，群）。映射 \((n, x) \mapsto x^{n}\) 把 \(N \times E\) （相应地 \(Z \times E\) ）映入 E ，由等式 \(x^{m+n} = x^{m}x^{n}\) （以加法作为 N 上的律），它关于第一个变量是分配的。若 E 是交换的，则该映射由等式 \((xy)^{n} = x^{n}y^{n}\) 关于第二个变量是分配的。

命题 1. 设 \(\mathbf{E}_1, \mathbf{E}_2, \ldots, \mathbf{E}_n\) 和 \(\mathbf{F}\) 为加法记法书写的交换幺半群，并设 \(u\) 为 \(\mathbf{E}_1 \times \cdots \times \mathbf{E}_n\) 到 \(\mathbf{F}\) 中的一个映射，它关于所有变量都是分配的。对于 \(i = 1, 2, \ldots, n\) ，设 \(\mathbf{L}_i\) 为非空有限集，且 \((x_{i,\lambda})_{\lambda \in \mathbf{L}_i}\) 为 \(\mathbf{E}_i\) 中元素的族。设 \(y_i = \sum_{\lambda \in \mathbf{L}_i} x_{i,\lambda}\) ， \(i = 1, 2, \ldots, n\) 。则

\[
u (y _ {1}, \dots , y _ {n}) = \sum_ {\alpha} u (x _ {1, \alpha_ {1}}, \dots , x _ {n, \alpha_ {n}})\tag{2}
\]

对所有序列求和 \(\alpha = (\alpha_{1},\ldots ,\alpha_{n})\) 的序列 \(\mathbf{L}_1\times \dots \times \mathbf{L}_n\)

我们通过对 \(n\) 进行归纳论证，基础情形 \(n = 1\) 由§1第2节的公式（2）得出。出自同一参考文献。

\[
u (y _ {1}, \dots , y _ {n - 1}, y _ {n}) = \sum_ {\alpha_ {n} \in L _ {n}} u (y _ {1}, \dots , y _ {n - 1}, x _ {n, \alpha_ {n}})\tag{3}
\]

对于 \(y_{n} = \sum_{\alpha_{n} \in \mathbf{L}_{n}} x_{n, \alpha_{n}}\) 且映射 \(z \mapsto u(y_{1}, \ldots, y_{n-1}, z)\) 于 \(\mathbf{E}_{n}\) 到 \(F\) 是一个原群同态。将归纳假设应用于分配映射 \((z_{1}, \ldots, z_{n-1}) \mapsto u(z_{1}, \ldots, z_{n-1}, x_{n, \alpha_{n}})\) 从 \(\mathbf{E}_{1} \times \cdots \times \mathbf{E}_{n-1}\) 到 \(F\) ,

\[
u (y _ {1}, \dots , y _ {n - 1}, x _ {n, \alpha_ {n}}) = \sum_ {\alpha_ {1}, \dots , \alpha_ {n - 1}} u (x _ {1, \alpha_ {1}}, \dots , x _ {n - 1, \alpha_ {n - 1}}, x _ {n, \alpha_ {n}}),\tag{4}
\]

求和是对这些序列进行的 \((\alpha_{1},\ldots,\alpha_{n-1})\) 的序列 \(M=L_{1}\times\cdots\times L_{n-1}\) 的一个子群。现在 \(L_{1}\times\cdots\times L_{n}=M\times L_{n}\) ；记

\[
t _ {\alpha_ {1}, \dots , \alpha_ {n}} = u (x _ {1, \alpha_ {1}}, \dots , x _ {n, \alpha_ {n}}),
\]

我们有

\[
\sum_ {\alpha_ {1}, \dots , \alpha_ {n}} t _ {\alpha_ {1}, \dots , \alpha_ {n}} = \sum_ {\alpha_ {n}} \left(\sum_ {\alpha_ {1}, \dots , \alpha_ {n}} t _ {\alpha_ {1}, \dots , \alpha_ {n - 1}, \alpha_ {n}}\right)\tag{5}
\]

由§1第5节的公式(7)，(2)立即由(3)、(4)和(5)得出。

注记。若 \(u(a_{1},\ldots ,a_{i - 1},0,a_{i + 1},\ldots ,a_{n}) = 0\) 对于 \(i = 1,2,\dots,n\) 和 \(a_{j}\in \mathbf{E}_{j}\) ( \(j\neq i\) )，则公式(2)对有限支撑的族 \((x_{i,\lambda})_{\lambda \in \mathbf{L}_i}\) 仍然成立。

定义5的一个特例是，u为与集合的作用相关联的作用法则的情形。 \(\Omega\) 在岩浆 E 上。如果 u 关于第二个变量是分配的，则也称该作用为 \(\Omega\) 在岩浆 E 上是分配的。换言之：

定义6. 一个作用 \(\alpha \mapsto f_{\alpha}\) 一个集合的 \(\Omega\) 在岩浆 E 上，若对所有 \(\alpha \in \Omega\) ，映射 \(f_{\alpha}\) 是岩浆 E 的一个自同态。

上定义了一个处处有定义的合成律。若 \(\top\) 表示岩浆 E 的运算，而 \(\bot\) 是与 \(\Omega\) 在 E 上的作用相伴的作用运算，后者的分配性则由公式

\[
\alpha \perp (x \top y) = (\alpha \perp x) \top (\alpha \perp y) \quad (\text { for   } \alpha \in \Omega \text {   and   } x, y \in E).\tag{6}
\]

表达。借用语言上的习惯，也称运算 \(\perp\) 关于运算 \(\top\).

是分配的（或右分配的）。§1 第 2 号中的公式 (2) 表明，此时，对 E 中元素的每个有序序列 \((x_{\lambda})_{\lambda \in L}\) 以及每个 \(\alpha \in \Omega\) ,

\[
\alpha \perp \left(\mathsf {T} _ {\lambda \in L} x _ {\lambda}\right) = \mathsf {T} _ {\lambda \in L} (\alpha \perp x _ {\lambda}).\tag{7}
\]

如果一个作用 \(\alpha\mapsto f_{\alpha}\) 是可分配的，并且 E 上的等价关系 R 与 E 的合成律及该作用相容 \(\alpha\mapsto f_{\alpha}\) ，则 E/R 上的商作用是可分配的。

当 E 上的律用乘法书写时，我们常使用指数记号 \(x^{\alpha}\) 来表示关于该乘法可分配的作用律，因此可分配性由恒等式 \((xy)^{\alpha} = x^{\alpha}y^{\alpha}\) 表达。若 E 上的律用加法书写，我们常使用左（相应地，右）乘法记号 \(\alpha.x\) （相应地 \(x.\alpha\) ）来表示关于该加法可分配的作用律，其可分配性由恒等式

\[
\alpha (x + y) = \alpha x + \alpha y \quad (\text { resp. } (x + y) \alpha = x \alpha + y \alpha).
\]

表达。我们也可以考虑 \(\Omega\) 具有一个内律，记为 \(\overline{\top}\) ，且作用律关于第一个变量可分配，这意味着

\[
(\alpha \overline{\top} \beta) \perp x = (\alpha \perp x) \top (\beta \perp x)\tag{8}
\]

为所有 \(\alpha, \beta \in \Omega\) 和 \(x \in \mathbf{E}\) 。于是，由 §1 第 2 号的公式 (2)

\[
\left(\overline {{\mathsf {T}}} _ {\lambda \in L} \alpha_ {\lambda}\right) \perp x = \mathsf {T} _ {\lambda \in L} (\alpha_ {\lambda} \perp x)\tag{9}
\]

对每一个有序序列 \((\alpha_{\lambda})_{\lambda \in L}\) 的 \(\Omega\) 以及所有 \(x\in E\) .

\subsection{一种内运算对另一种内运算的分配律}

定义7。设 \(\top\) 和 \(\bot\) 是集合 E 上的两个内运算。该运算 \(\bot\) 被称为关于该法则分配的 \(\top\) 中正规，若

\begin{enumerate}
\def\labelenumi{(\arabic{enumi})}
\setcounter{enumi}{9}
\tightlist
\item
\end{enumerate}

\[
x \perp (y \top z) = (x \perp y) \top (x \perp z)\tag{11}
\]

\[
(x \top y) \perp z = (x \perp z) \top (y \perp z)
\]

为所有 \(x, y, z\) 在 \(\mathbf{E}\) .

注意，若该定律成立，则(10)与(11)等价。 \(\perp\) 是可交换的。一般而言，其中一个运算律以加法形式书写，另一个以乘法形式书写；若乘法对加法满足分配律，则：

\begin{enumerate}
\def\labelenumi{(\arabic{enumi})}
\setcounter{enumi}{11}
\tightlist
\item
\end{enumerate}

\[
x. (y + z) = x. y + x. z\tag{13}
\]

\[
(x + y). z = x. z + y. z
\]

示例。(1) 在集合中 \(\mathfrak{P}(\mathbf{E})\) 的子集集合 \(\mathbf{E}\) ，每个内部运算 \(\cap\) 和 \(\cup\) 对自身和另一个运算都是分配的。这由形如

\[
\begin{array}{l} \mathbf {A} \cap (\mathbf {B} \cup \mathbf {C}) = (\mathbf {A} \cap \mathbf {B}) \cup (\mathbf {A} \cap \mathbf {C}) \\ \mathbf {A} \cup (\mathbf {B} \cap \mathbf {C}) = (\mathbf {A} \cup \mathbf {B}) \cap (\mathbf {A} \cup \mathbf {C}). \end{array}
\]

\begin{enumerate}
\def\labelenumi{(\arabic{enumi})}
\setcounter{enumi}{1}
\item
  在 Z 中（更一般地，在任何全序集中），sup 和 inf 这两个运算中的每一个都对另一个运算以及自身是分配的。
\item
  在 \(\mathbf{Z}\) （*更一般地，在任何环中*）乘法对加法是分配的。
\item
  在 N 中，加法和乘法对 sup 和 inf 这两个运算都是分配的。
\end{enumerate}

\section{群与带算子的群}

\subsection{群}

回顾以下定义（§2，第3号，定义6）。

定义1. 一个具有结合律的合成法则、含有单位元且其中每个元素都可逆的集合，称为群。

换言之，群是其中每个元素都可逆的幺半群（§2，第1号，定义1）。集合上确定其群结构的合成律称为群律。若G和H是两个群，则从G到H的岩浆同态也称为群同态。这样的同态f将单位元映至单位元；事实上，设e（分别为e'）为G（分别为H）的单位元；将G和H的群律写成乘法形式，有e·e = e，从而 \(f(e).f(e) = f(e)\) 且乘以 \(f(e)^{-1}, f(e) = e'\) 因此 f 是含幺的。于是由第 2 节第 3 条可知 \(f(x^{-1}) = f(x)^{-1}\) 为所有 \(x \in G\) .

例。在任何幺半群 E 中，可逆元素集在 E 上诱导的结构下构成一个群。特别地，集合 F 到自身的双射映射集（或 F 的置换集）在运算 \((f, g) \mapsto f \circ g\) 下构成一个群，称为集合 F 的对称群，记作 \(S_{F}\) .

在本段中，除非另有说明，群的合成法则将始终以乘法形式书写，且e将表示该群法则的单位元。

群G称为有限群，如果G的底集是有限的；否则称为无限群；群的基数称为群的阶。

若G上的一个合成法则确定了G上的群结构，其对偶法则亦然。将群G映到自身的映射，将每个元素对应到 \(x \in G\) x 的逆是 G 到其相反群的一个同构。 \(\S 2\) ，第3号，命题4）。

按照我们的一般约定（《集合论》第二卷，第3节，第1号），我们将用 \(\mathbf{A}^{-1}\) 表示映射下G的子集A的像。 \(x \mapsto x^{-1}\) 但重要的是要注意，尽管记号类似， \(\mathbf{A}^{-1}\) 绝不是A在G的子集之间的合成法则下的逆元素（回忆一下，XY是满足 \((\mathbf{X}, \mathbf{Y}) \mapsto \mathbf{XY}\) 的xy的集合）：该法则下的单位元是 \(x \in \mathbf{X}, y \in \mathbf{Y}\) ，而在此法则下 \(\{e\}\) 中仅有的可逆元素是由单个元素组成的集合A（这样的A，当然，其逆元素确实是 \(\mathfrak{P}(\mathbf{G})\) ）。单位元 \(\mathbf{A}^{-1}\) 。A称为 \((\mathbf{AB})^{-1} = \mathbf{B}^{-1}\mathbf{A}^{-1}\) 对……成立 \(\mathbf{A} \subset \mathbf{G}, \mathbf{B} \subset \mathbf{G}\) 的对称子集， \(\mathbf{G}\) 中正规，若 \(\mathbf{A} = \mathbf{A}^{-1}\) 的分拆集合的双射。对于所有 \(\mathbf{A} \subset \mathbf{G}, \mathbf{A} \cup \mathbf{A}^{-1}, \mathbf{A} \cap \mathbf{A}^{-1}\) 和 \(\mathbf{AA}^{-1}\) 是对称的。

\subsection{带算子的群}

定义2. 设 \(\Omega\) 为一个集合。一个群 G 连同 \(\Omega\) 在 G 上的一个作用，若该作用关于群律是分配的，则称 G 为以 \(\Omega\) .

为算子群的带算子群。在下文中， \(x^{\alpha}\) 将表示 \(\alpha\in\Omega\) 和 \(x\in G\) 的复合。分配性则由恒等式 \((xy)^{\alpha}=x^{\alpha}y^{\alpha}\) .

表达。在一个带算子群 G 中，每个算子定义了底层群结构的一个自同态；这些自同态有时被称为带算子群 G 的位似变换。

一个带算子群 G 称为交换的（或阿贝尔的），如果其群律是交换的。

在下文中，一个群 G 将被等同于以 \(\varnothing\) 为算子群的带算子群，这是通过赋予 G 以 \(\varnothing\) 在 G 上的唯一作用得到的。这使我们能够将群视为带算子群的特殊情形，并将我们即将陈述的关于后者的定义和结果应用于群。

例。在一个交换群 G 中，以乘法记号书写， \((xy)^{n} = x^{n}y^{n}\) 为所有 \(n \in Z\) （§2，第8号，方程（1））；因此 Z 在 G 上的作用 \(n \mapsto (x \mapsto x^{n})\) 连同群律一起，定义了 G 上带算子群的结构。

定义3. 设 G 和 G' 是以 Ω 为算子群的带算子群。从 G 到 G' 的带算子群同态是从群 G 到群 G' 的一个同态，使得

\[
f (x ^ {\alpha}) = (f (x)) ^ {\alpha}
\]

为所有 \(\alpha \in \Omega\) 以及所有 \(x \in G\) .

成立。带算子群 G 的自同态是群 G 的一个自同态，并且它与 G 的所有位似变换可交换。

由于带算子群 G 的两个位似变换不一定可交换，G 的一个位似变换一般不一定是带算子群 G 的自同态。

带算子群的恒等映射是带算子群的同态；两个带算子群同态的复合也是一个带算子群同态。一个映射是带算子群的同构，当且仅当它是带算子群的既单且满的同态，并且此时逆映射也是带算子群的同构。

更一般地，设 G（相应地 G'）是以 Ω（相应地 Ω'）为算子群的带算子群。设 φ 是从 Ω 到 Ω' 的一个映射。从 G 到 G' 的一个 φ-同态是从群 G 到群 G' 的一个同态，使得

\[
f (x ^ {\alpha}) = (f (x)) ^ {\phi (\alpha)}
\]

为所有 \(\alpha \in \Omega\) 以及所有 \(x \in G\) .

在本段的其余部分，我们将给定一个集合 \(\Omega\) 。除非另有说明，所考虑的带算子群将以 \(\Omega\) 作为算子集合。

\subsection{子群}

定义4. 设 G 是一个带算子群。G 的一个稳定子群是 G 的一个子集 H，它具有以下性质：

\begin{enumerate}
\def\labelenumi{(\roman{enumi})}
\item
  \(e\in \mathbf{H};\)
\item
  \(x, y \in \mathbf{H}\) 蕴含 \(xy \in \mathbf{H}\) ;
\item
  \(x \in \mathbf{H}\) 蕴含 \(x^{-1} \in \mathbf{H}\) ;
\item
  \(x \in \mathbf{H}\) 和 \(\alpha \in \Omega\) 蕴含 \(x^{\alpha} \in \mathbf{H}\) .
\end{enumerate}

如果 H 是 G 的一个稳定子群，则 H 上由 G 的带算子群结构诱导的结构是带算子群的结构，并且 H 到 G 的典范单射是带算子群的同态。

设 G 是一个群。G 的一个稳定子群，连同 \(\varnothing\) 的作用（第2号），即满足定义4的条件（i）、（ii）、（iii）的 G 的子集，称为 G 的子群。当我们谈到一个带算子群的子群时，我们总是指 G 的底层群的子群。带算子群 G 的一个子群不一定是 G 的稳定子群。

例（1）。设 \(\Sigma\) 是结构的一个种（集合论，IV，§1，第4号），而S是种的一个结构 \(\Sigma\) 在集合E上（同上）。S的自同构集合是……的一个子群 \(\mathfrak{S}_{\mathbb{E}}\) .

命题1. 设G是一个带算子的群，H是G的一个在G的位似变换下稳定的子集。以下条件是等价的：

\begin{enumerate}
\def\labelenumi{(\alph{enumi})}
\item
  H 是 G 的一个稳定子群。
\item
  H 是非空的，且关系 \(x \in \mathbf{H}, y \in \mathbf{H}\) 蕴含 \(xy \in \mathbf{H}\) 和 \(x^{-1} \in \mathbf{H}\) .
\item
  H 是非空的，且关系 \(x \in \mathbf{H}\) , \(y \in \mathbf{H}\) 蕴含 \(xy^{-1} \in \mathbf{H}\) .
\item
  H在G上的运算下是稳定的，且由G上的运算在H上诱导出的合成运算构成一个群运算。
\end{enumerate}

显然(a)蕴含(b)。我们证明(b)蕴含(a)。只需证明H包含G的单位元即可。由于子集H非空，设 \(x \in H\) 。则 \(x^{-1} \in H\) 和 \(e = xx^{-1} \in H\) 显然(b)蕴含(c)。我们证明(c)蕴含(b)。首先，由于H非空，它包含一个元素x。因此 \(xx^{-1} = e\) 是 H 的一个元素。对于 H 中的每一个元素 \(x\) ， \(x^{-1} = ex^{-1}\) 属于 H；因此这些关系 \(x \in H\) , \(y \in H\) 蕴含 \(x(y^{-1})^{-1} = xy \in H\) 。显然 (a) 蕴含 (d)。我们证明 (d) 蕴含 (a)：H 到 G 的典范嵌入是一个群同态；因此 \(e \in H\) 以及关系 \(x \in H\) 蕴含 \(x^{-1} \in H\) （编号 1）。

注：(1) 类似地可以证明，条件(b)等价于条件

(c') H ≠ ∅ 且关系 \(x \in H\) 和 \(y \in H\) 蕴含 \(y^{-1}x \in H\) .

\begin{enumerate}
\def\labelenumi{(\arabic{enumi})}
\setcounter{enumi}{1}
\tightlist
\item
  对于G的每个子群H，有以下关系：
\end{enumerate}

\[
\mathrm{H.H} = \mathrm{H} \quad \text { and } \quad \mathrm{H} ^ {- 1} = \mathrm{H}.\tag{1}
\]

对于 H.H ⊂ H 且 H \(^{-1}\) 由 (b) 知 ⊂ H。由于 e \ensuremath{\in} H，H.H ⊃ e.H = H，且取逆将包含关系 H \(^{-1}\) ⊂ H 到 H ⊂ H \(^{-1}\) ，由此得到公式 (1)。

若H是G的稳定子群，且K是H的稳定子群，显然K也是G的稳定子群。

集合 \(\{e\}\) 是G的最小稳定子群。G的一族稳定子群的交集仍是稳定子群。因此，存在包含G的给定子集X的最小稳定子群H；它被称为由X生成的稳定子群，而X被称为H的一个生成系（或生成集）。

命题2。设X是带算子群G的一个非空子集，且 \(\hat{X}\) 在作用下的稳定子集 \(\Omega\) 在由X生成的G上。由X生成的稳定子群是由集合Y = \(\hat{X} \cup \hat{X}^{-1}\) .

后者子集 Z 是所有项均为某集合元素的有限序列的复合构成的集合。 \(\hat{X}\) 或元素的逆元 \(\hat{X}\) ：这种复合的逆元是具有相同形式的复合（§2，第3号，命题5的推论1），且Z在 \(\Omega\) 正如通过将§3第4号命题1应用于G的位似变换所见，因此（命题1）Z是G的一个稳定子群。反之，每个包含X的稳定子群显然包含Y，从而包含Z。

推论1. 设G是一个带算子的群，X是G中在Ω作用下稳定的子集。由X生成的子群与由X生成的稳定子群相同。

推论2. 设G是一个群，X是G中由两两可置换元素组成的子集。由X生成的G的子群是交换的。

集合 \(Y = X \cup X^{-1}\) 由两两可置换的元素组成（§2第3号命题5），且在由Y生成的稳定子集上诱导的运算是交换的（§1第5号推论2）。

如果G是一个带算子的群，由G中两两可置换元素组成的子集所生成的稳定子群不一定是交换的。

推论3。设 \(f: \mathbf{G} \to \mathbf{G}'\) 设为一个带算子的群同态，且 \(\mathbf{X}\) 的一个子集 \(\mathbf{G}\) 。在 \(f\) 稳定子群的一个子集 \(\mathbf{G}\) 生成的 \(\mathbf{X}\) 是 \(\mathbf{G}'\) 生成的 \(f(\mathbf{X})\) .

设 \(\mathbf{X}' = f(\mathbf{X})\) 。则 \(\hat{\mathbf{X}}' = f(\hat{\mathbf{X}})\) 和 \(\mathbf{X}^{\prime -1} = f(\mathbf{X}^{-1})\) 。因此

\[
f (\hat {\mathbf {X}} \cup \hat {\mathbf {X}} ^ {- 1}) = \hat {\mathbf {X}} ^ {\prime} \cup \hat {\mathbf {X}} ^ {- 1}.
\]

推论随后由§1第4号命题1得出。

例（2）。设 G 是一个群，x 是 G 的一个元素。由 \(\{x\}\) （更简单地称为由 x 生成的子群）是集合 \(x^{n}, n \in Z\) 在G上的律作用下生成的稳定子集 \(\{x\}\) 是 \(x^{n}\) （其中 \(n \in N^{*}\) 这两个集合一般而言是不同的。

因此，在加法群 \(\mathbf{Z}\) 中，由元素 \(x\) 是集合 \(x.\mathbf{Z}\) 于 \(xn\) , \(n \in \mathbf{Z}\) ，以及由 \(x\) 是 \(xn\) , \(n \in \mathbf{N}^*\) 这两个集合在以下情况下总是不同的 \(x \neq 0\) .

右向族稳定子群之并显然是G的稳定子群。由此可知，若P是G的子集且H是G中不与P相交的稳定子群，则包含H且不与P相交的G的稳定子群之集，按包含关系排序，是归纳的（集合论，III，§2，第4节）。应用Zorn引理（集合论，III，§2，第4节），我们得到如下结果：

命题3. 设G是一个带算子的群，P是G的一个子集，H是G的一个稳定子群且与P不相交。G中包含H且与P不相交的稳定子群的集合具有一个极大元。

\subsection{商群}

定理1. 设R是带算子群G上的一个等价关系；若R与G上的群律左（相应地，右）相容（§3，第3号）且与Ω的作用相容，则e的等价类为G的一个稳定子群H，且关系R等价于 \(x^{-1}y \in H\) （相应地 \(yx^{-1} \in H\) ）。反之，若H为G的一个稳定子群，则关系 \(x^{-1}y \in H\) （相应地 \(yx^{-1} \in H\) 是一个等价关系，它在左侧（相应地，右侧）与 G 上的群运算相容，并且与 Ω 的作用相容，在此关系下 H 是 e 的等价类。

我们将注意力限制在关系 R 与 G 上的运算左相容的情形（右相容关系的情形可通过将 G 上的运算替换为相反运算来得到）。该关系 \(y \equiv x \pmod{R}\) 等价于 \(x^{-1}y \equiv e \pmod{R}\) ，对于 \(y \equiv x\) 蕴含 \(x^{-1}y \equiv x^{-1}x = e\) 反之亦然 \(x^{-1}y \equiv e\) 蕴含 \(y = x(x^{-1}y) \equiv x\) 。若 H 表示 e 的等价类，则关系 R 等价于 \(x^{-1}y \in H\) 。我们证明 H 是 G 的一个稳定子群。对于每个算子 \(\alpha\) ，关系 \(x \equiv e\) 蕴含 \(x^{\alpha} \equiv e^{\alpha} = e\) ，因此 \(H^{\alpha} \subset H\) 且 H 在 \(\Omega\) 的作用下是稳定的。只需建立（命题 1） \(x \in H\) 和 \(y \in H\) 蕴含 \(x^{-1}y \in H\) ，即 \(x \equiv e\) 和 \(y \equiv e\) 蕴含 \(x \equiv y\) ，这是 R 的传递性的一个推论。

反之，设H是G的一个稳定子群；该关系 \(x^{-1}y \in H\) 是自反的，因为 \(x^{-1}x = e \in H\) ；它是对称的，因为 \(x^{-1}y \in H\) 蕴含 \(y^{-1}x = (x^{-1}y)^{-1} \in H\) 它是传递的，因为 \(x^{-1}y \in H\) 和 \(y^{-1}z \in H\) 蕴含 \(x^{-1}z = (x^{-1}y)(y^{-1}z) \in H\) 它与G上的合成律左相容，因为 \(x^{-1}y = (zx)^{-1}(zy)\) 为所有 \(z \in G\) 最后，对于每个算子 \(\alpha\) ，关系 \(y \in xH\) 蕴含 \(y^{\alpha} \in x^{\alpha}H^{\alpha} \subset x^{\alpha}H\) 因此等价关系 \(x^{-1}y \in H\) 与 \(\Omega\) 在G上的作用相容。

设G是一个群，H是G的一个子群；关系 \(x^{-1}y \in H\) （相应地 \(yx^{-1} \in H\) )也写成等价形式 \(y \in xH\) （相应地 \(y \in Hx\) ). 因此，G 的每个子群 H 都在 G 上定义了两种等价关系，即 \(y \in xH\) 和 \(y \in Hx\) 在这些关系下的等价类分别是集合xH，称为H的左陪集（或模H的左陪集），以及集合Hx，称为H的右陪集（或模H的右陪集）。通过关于这些关系对子集A ⊂ G进行饱和（集合论，II，§6，第4节），我们分别得到集合AH和HA。映射 \(x \mapsto x^{-1}\) 将模H的左陪集变换为模H的右陪集，反之亦然。

左陪集（模 H）的集合的基数称为子群 H 关于 G 的指数，记为 (G:H)；它也等于右陪集的集合的基数。

若G的子群K包含H，则K是H的左（或右）陪集的并集。由于K的左陪集由K经左平移得到，包含在K的一个左陪集内的H的左陪集所成集合的基数与后者无关。因此（集合论，III，§5，第8节，命题9）：

命题4。设H和K是群G的两个子群，且H ⊂ K。则

\[
(\mathbf {G}: \mathbf {H}) = (\mathbf {G}: \mathbf {K}) (\mathbf {K}: \mathbf {H}).\tag{2}
\]

推论。若G是一个阶为g的有限群，且H是G的一个阶为h的子群，则

\[
h. (\mathbf {G}: \mathbf {H}) = g\tag{3}
\]

（特别地，H 的阶和指数都是 G 的阶的因数。）

定理1使我们能够确定与带算子群G上的运算相容的等价关系：若R是这样一个关系，则它既与G上的群运算左相容又右相容，并且与算子的作用相容。 \(\Omega\) 因此，若H是e（模R）的类，则H是一个稳定子群，使得关系 \(y \in xH\) 和 \(y \in Hx\) 是等价的（因为两者都等价于R）；因此 \(xH = Hx\) 为所有 \(x \in G\) 。反之，若如此，则等价关系之一 \(y \in xH\) , \(y \in Hx\) 与群律相容，因为它与该律既左相容又右相容（§3，第4节），并且与 \(\Omega\) 的作用相容。由于方程 \(xH = Hx\) 等价于 \(xHx^{-1} = H\) ，我们作出如下定义：

定义5. 设G是一个带算子的群。G的一个稳定子群H称为G的一个正规（或不变）稳定子群，如果 \(xHx^{-1} = H\) 为所有 \(x \in G\) .

上定义了一个处处有定义的合成律。若 \(\Omega = \varnothing\) ，G的一个正规稳定子群称为G的一个正规（或不变）子群。在交换群中，每个子群都是正规的。

要验证一个稳定子群H是正规的，只需证明 \(xHx^{-1} \subset H\) 为所有 \(x \in G\) ；因为若如此，则 \(x^{-1}Hx \subset H\) 为所有 \(x \in G\) ，即 \(H \subset xHx^{-1}\) ，从而 \(H = xHx^{-1}\) .

设H是G的一个正规稳定子群，R是由H定义的等价关系 \(y \in xH\) ；在商集G/R上，内运算（即群G的运算关于R的商）是结合的；e的类是此商运算的单位元；G中两个互逆元素的类在商运算下互为逆元，并且 \(\Omega\) 的作用（即 \(\Omega\) 在G上的作用关于R的商）关于G/R上的内运算是分配的（ \(\S 3\) ，第5节）。因此，总结所得结果：

定理2. 设G是一个带算子的群。对于G上的一个等价关系R，要使它与群运算及 \(\Omega\) 的作用相容，必要且充分条件是它具有形式 \(x^{-1}y \in H\) ，其中H是G的一个正规稳定子群（对于这样的子群，关系 \(x^{-1}y \in H\) 还等价于 \(yx^{-1} \in H\) ）。G/R上的合成运算（即G上运算的商）以及 \(\Omega\) 在G/R上的作用（即 \(\Omega\) 在G上的作用关于这样的关系R的商）赋予G/R带算子的群的结构，称为商结构，并且到商集的典范映射是带算子的群的同态。

定义6. 带算子的群G关于由G的正规子群H定义的等价关系的商，连同商结构，称为G关于H的带算子的商群，记为G/H。典范映射G → G/H称为典范同态。

设G是一个群，H是G的一个正规子群。商G/H连同其群结构称为G关于H的商群。对于从G/H到某个带算子的群的映射，要使它是带算子的群的同态，必要且充分条件是它与从G到G/H的典范映射的复合是同态：这证明了“商群”这一名称的合理性（集合论，IV，§ 2，第6节）。

由G的正规稳定子群定义的等价关系记为 \(x \equiv y \pmod{H}\) 或 \(x \equiv y(H)\) .

命题5. 设 \(f: G \to G'\) 是带算子的群的同态，H和H'分别是G和G'的正规稳定子群，使得 \(f(H) \subset H'\) 该映射 \(f\) 与由H和H'定义的等价关系相容。设 \(\pi: G \to G/H\) 和 \(\pi': G' \to G'/H'\) 是典范同态。由 \(\bar{f}: G/H \to G'/H'\) 通过取商导出的映射 \(f\) 是一个同态。

上定义了一个处处有定义的合成律。若 \(x \equiv y \pmod{H}\) ，则 \(x^{-1}y \in H\) 的自同态，因此

\[
f (x) ^ {- 1} f (y) = f \left(x ^ {- 1}\right) f (y) = f \left(x ^ {- 1} y\right) \in f (\mathrm{H}) \subset \mathrm{H} ^ {\prime}
\]

，因此 \(f(x) \equiv f(y) \pmod{H'}\) 第二个断言由商律的泛性质（§1，第6节）得出。

注记。(1) 若A是群G的任意子集，H是G的正规子群，则AH = HA；该集合通过对A关于该关系进行饱和而得到。 \(x \equiv y \pmod{H}\) .

\begin{enumerate}
\def\labelenumi{(\arabic{enumi})}
\setcounter{enumi}{1}
\tightlist
\item
  若H是G的有限指数的正规子群，则商群G/H是阶为(G:H)的有限群。
\end{enumerate}

注意，若H是群G的正规子群，且K是H的正规子群，则K不一定是G的正规子群（I，§5，习题10）。

设G是一个带算子的群。G的所有正规稳定子群的任意族的交是正规稳定子群。因此，对于G的每个子集X，存在包含X的最小正规稳定子群，称为由X生成的正规稳定子群。

在具有算子群 G 中，稳定子群 G 和 \(\{e\}\) 是正规的。

定义 7. 具有算子群 G 称为单的，如果 \(G \neq \{e\}\) 并且 G 中不存在除 G 和 \(\{e\}\) .

\subsection{同态的分解}

命题6. 设G是一个带算子的群，G'是一个带有Ω作用的原群，作用以指数形式书写。设f: G → G'是原群G到原群G'的一个同态，使得对所有α \ensuremath{\in} Ω及所有x \ensuremath{\in} G， \(f(x^{\alpha}) = f(x)^{\alpha}\) 。则 \(f(\mathbf{G})\) 是G'在G'上的运算及Ω作用下的稳定子集；该集合 \(f(\mathbf{G})\) 在诱导运算下是一个带算子的群，且该映射 \(x \mapsto f(x)\) 是一个群同态，其核为 G 的中心，其像为 \(f(\mathbf{G})\) 是带算子群的同态。

根据§1第4号命题1， \(f(\mathbf{G})\) 是……的稳定子集 \(\mathbf{G}'\) 根据内部法律规定， \(\mathbf{G}'\) 。对于每个元素 \(x \in \mathbf{G}\) 以及对于每个算子 \(\alpha\) , \(f(x)^{\alpha} = f(x^{\alpha}) \in f(\mathbf{G})\) ，从而 \(f(\mathbf{G})\) 在 \(\Omega\) 在 \(\mathbf{G}'\) 的作用下是稳定的。将 \(\mathbf{G}'\) 的内部运算写成乘法形式，

\[
\begin{array}{r l} (f (x) f (y)) f (z) & = f (x y) f (z) = f ((x y) z) = f (x (y z)) = f (x) f (y z) \\ & \qquad \qquad \qquad \qquad \qquad \qquad \qquad \qquad \qquad \qquad = f (x) (f (y) f (z)) \end{array}
\]

对于所有元素 \(x, y, z\) 在 \(\mathbf{G}\) ；因此诱导的定律在 \(f(\mathbf{G})\) 是结合的。

设 \(e\) 是 \(G\) 的单位元。其像 \(f(e)\) 是 \(f(G)\) 的单位元（§ 2，第 1 号）。的每个元素 \(f(G)\) 在 \(f(G)\) (§ 2，第 3 号)。因此，在……上诱导的法则 \(f(G)\) 根据内部法律规定， \(G'\) 是一个群法则。对于所有元素 \(x\) 和 \(y\) 在 \(G\) 以及每个算子 \(\alpha\) ,

\[
(f (x) f (y)) ^ {\alpha} = (f (x y)) ^ {\alpha} = f ((x y) ^ {\alpha}) = f \left(x ^ {\alpha} y ^ {\alpha}\right) = f \left(x ^ {\alpha}\right) f \left(y ^ {\alpha}\right) = (f (x)) ^ {\alpha} (f (y)) ^ {\alpha}
\]

这表明……的作用 \(\Omega\) 对群运算具有分配性 \(f(\mathrm{G})\) 的单位同态。因此 \(f(\mathrm{G})\) 在诱导运算下是一个带算子的群，并且显然该映射 \(x \mapsto f(x)\) 是带算子群的同态。

定义8. 设 \(f: \mathbf{G} \to \mathbf{G}'\) 设为单位元的逆像。 \(\mathbf{G}'\) 被称为 \(f\) .

的核 \(f\) 通常记作 \(\operatorname{Ker}(f)\) 以及图像 \(f(\mathbf{G})\) 于 \(f\) 有时用……表示 \(\operatorname{Im}(f)\) .

定理 3. 设 \(f: \mathbf{G} \to \mathbf{G}'\) 设为一个带算子的群同态。

\begin{enumerate}
\def\labelenumi{(\alph{enumi})}
\item
  \(\operatorname{Ker}(f)\) 是 \(\mathbf{G}\) ;
\item
  \(\operatorname{Im}(f)\) 是 \(\mathbf{G}'\) ;
\item
  映射 \(f\) 与定义在……上的等价关系相容 \(\mathbf{G}\) 中条件的扩张 \(\operatorname{Ker}(f)\)
\item
  映射 \(\tilde{f}\colon \mathbf{G} / \mathbf{Ker}(f)\to \operatorname {Im}(f)\) 通过取商导出的映射 \(f\) 通过取商，这是带算子的群的同构；
\item
  \(f = \iota \circ \tilde{f} \circ \pi\) 的x和y，其中 \(\iota\) 是以下映射的典范嵌入 \(\operatorname{Im}(f)\) 到 \(G'\) 和 \(\pi\) 是 \(G\) 到 \(G / \operatorname{Ker}(f)\) .
\end{enumerate}

的典范同态。断言(b)由命题6得出。等价关系 \(f(x) = f(y)\) 在 \(G\) 与 \(G\) 上的带算子群结构相容。由定理2（第4节），它因此具有形式 \(y \in xH\) 的x和y，其中 \(H\) 是 \(G\) 和 \(H\) 是单位元的类，因此 \(H = \operatorname{Ker}(f)\) 。于是断言(a)、(c)和(d)成立。断言(e)是显然的（集合论，II，§6，第5节）。

\subsection{商群的子群}

命题7. 设G和H是两个带算子的群，f是G到H的一个同态，N是f的核。

\begin{enumerate}
\def\labelenumi{(\alph{enumi})}
\item
  设 \(\mathbf{H}'\) 是 \(\mathbf{H}\) 的一个稳定子群。原像 \(\mathbf{G}' = f^{-1}(\mathbf{H}')\) 是 \(\mathbf{G}\) 和 \(\mathbf{G}'\) 在 \(\mathbf{G}\) 中正规，若 \(\mathbf{H}'\) 在 \(\mathbf{H}\) 。进一步， \(\mathbf{N}\) 是由 \(\mathbf{G}'\) 的一个正规子群。若 \(f\) 是满射，则 \(\mathbf{H}' = f(\mathbf{G}')\) 和 \(f\) 在取商后定义了 \(\mathbf{G}' / \mathbf{N}\) 到 \(\mathbf{H}'\) 的一个同构。
\item
  设 \(\mathbf{G}'\) 是 \(\mathbf{G}\) 。像 \(\mathbf{H}' = f(\mathbf{G}')\) 是 \(\mathbf{H}\) 和 \(f^{-1}(\mathbf{H}') = \mathbf{G}'\mathbf{N} = \mathbf{NG}'\) 。特别地， \(f^{-1}(\mathbf{H}') = \mathbf{G}'\) 当且仅当 \(\mathbf{N} \subset \mathbf{G}'\) 的一个正规子群。若 \(f\) 是满射且 \(\mathbf{G}'\) 在 \(\mathbf{G}\) ，则 \(\mathbf{H}'\) 在 \(\mathbf{H}\) .
\item
  设 \(x\) 和 \(y\) 的序列 \(G'\) 和 \(\alpha \in \Omega\) ；于是 \(f(x) \in H'\) 和 \(f(y) \in H'\) 的自同态，因此 \(f(xy^{-1}) = f(x)f(y)^{-1} \in \mathbf{H}'\) ，即 \(xy^{-1} \in \mathbf{G}'\) ；因此 \(\mathbf{G}'\) 是 \(\mathbf{G}\) 的一个子群。现在 \(f(x^{\alpha}) = f(x)^{\alpha} \in \mathbf{H}'\) 的自同态，因此 \(x^{\alpha} \in \mathbf{G}'\) ，从而 \(\mathbf{G}'\) 是稳定的。假设 \(\mathbf{H}'\) 在 \(\mathbf{H}\) ，并设 \(x \in \mathbf{G}'\) , \(y \in \mathbf{G}\) ；于是 \(f(x) \in \mathbf{H}'\) 和
\end{enumerate}

\[
f (y x y ^ {- 1}) = f (y) f (x) f (y) ^ {- 1} \in \mathrm{H} ^ {\prime}
\]

，因此 \(yxy^{-1} \in G'\) ；因此 \(G'\) 在 \(G\) 的分拆集合的双射。对于所有 \(n \in N\) , \(f(n) = e \in H'\) 的自同态，因此 \(N \subset G'\) ；由于 \(N\) 在 \(G\) ，它在 \(G'\) 中是正规的。最后，若 \(f\) 是满射的， \(f(f^{-1}(A)) = A\) 对每个子集 \(A\) 于 \(H\) 的自同态，因此 \(H' = f(G')\) 成立； \(f\) 到 \(G'\) 的限制是核为 \(f'\) 于 \(G'\) 到 \(H'\) 的同态 \(N\) ，因此 \(f'\) 在过渡到商后定义了 \(G'/N\) 到 \(H'\) .

\begin{enumerate}
\def\labelenumi{(\alph{enumi})}
\setcounter{enumi}{1}
\tightlist
\item
  设 \(a\) 和 \(b\) 的序列 \(\mathbf{H}'\) 和 \(\alpha\) 在 \(\Omega\) 的一个同构；存在 \(x, y\) 在 \(\mathbf{G}'\) 使得 \(a = f(x)\) 和 \(b = f(y)\) 的自同态，因此 \(ab^{-1} = f(xy^{-1}) \in \mathbf{H}'\) ，因此 \(\mathbf{H}'\) 是 \(\mathbf{H}\) ，它对 \(a^{\alpha} = f(x^{\alpha}) \in \mathbf{H}'\) 的一个映射。设 \(x \in \mathbf{G}\) ；于是 \(x \in f^{-1}(\mathbf{H}')\) 当且仅当 \(f(x) \in \mathbf{H}' = f(\mathbf{G}')\) ，即当且仅当存在 \(y\) 在 \(\mathbf{G}'\) 使得 \(f(x) = f(y)\) ；关系 \(f(x) = f(y)\) 等价于存在 \(n \in \mathbf{N}\) 使得 \(x = yn\) ；最后， \(x \in f^{-1}(\mathbf{H}')\) 等价于 \(x \in \mathbf{G}'\mathbf{N} = \mathbf{NG}'\) . 显然关系 \(\mathbf{G}' = \mathbf{G}'\mathbf{N}\) 等价于 \(\mathbf{G}' \supset \mathbf{N}\) . 最后假设 \(f\) 是满射且 \(\mathbf{G}'\) 在 \(\mathbf{G}\) ；令 \(a \in \mathbf{H}'\) 和 \(b \in \mathbf{H}\) 的一个同构；存在 \(x \in \mathbf{G}'\) 和 \(y \in \mathbf{G}\) 使得 \(a = f(x)\) 和 \(b = f(y)\) 的自同态，因此 \(bab^{-1} = f(yxy^{-1}) \in f(\mathbf{G}') = \mathbf{H}'\) 。因此 \(\mathbf{H}'\) 在 \(\mathbf{H}\) .
\end{enumerate}

推论1. 假设 \(f\) 是满射。设 \(\mathfrak{G}\) （相应地 \(\mathfrak{G}'\) ) 为稳定（相应地，正规稳定）子群构成的集合，这些子群包含 \(\mathbf{G}\) ，包含 \(\mathbf{N}\) 和 \(\mathfrak{H}\) （相应地 \(\mathfrak{H}'\) ) 为稳定（相应地，正规稳定）子群构成的集合，这些子群包含 \(\mathbf{H}\) ，这些集合按包含关系排序。映射 \(\mathbf{G}' \mapsto f(\mathbf{G}')\) 是一个有序集同构 \(\Phi: \mathfrak{G} \to \mathfrak{H}\) ；逆同构 \(\Psi: \mathfrak{H} \to \mathfrak{G}\) 是映射 \(\mathbf{H}' \mapsto f^{-1}(\mathbf{H}')\) 。进一步 \(\Phi\) 和 \(\Psi\) 诱导同构 \(\Phi': \mathfrak{G}' \to \mathfrak{H}'\) 和 \(\Psi': \mathfrak{H}' \to \mathfrak{G}'\) .

推论2. 设 \(f: G \to H\) 设 是带算子群的同态，N 为其核 \(f\) , \(G'\) G的一个稳定子群，且L是G的一个正规稳定子群 \(G'\) 则 LN, L. ( \(G' \cap N\) ) 且 \(f(L)\) 是 \(G'N\) , \(G'\) 和 \(f(G')\) 分别，以及三个带算子的商群 \(G'N/LN\) , \(G'/L\) . ( \(G' \cap N\) ) 且 \(f(G')/f(L)\) 是同构的。

设 \(H' = f(G')\) 和 \(f'\) 表示从 \(G'\) 到 \(H'\) 它与……一致 \(f\) 在 \(G'\) ；其核 \(f'\) 是 \(G' \cap N\) 和 \(f'(L) = f(L)\) 由命题7可知， \(f'(L)\) 是 \(H'\) 和

\[
f ^ {\prime - 1} (f ^ {\prime} (\mathbf {L})) = \mathbf {L}. (\mathbf {G} ^ {\prime} \cap \mathbf {N})
\]

是 \(\mathbf{G}'\) 的一个映射。设 \(\lambda\) 为典范同态 \(\mathbf{H}'\) 到 \(\mathbf{H}' / f'(\mathbf{L}) = f(\mathbf{G}') / f(\mathbf{L})\) ：作为 \(\lambda \circ f'\) 是满射且核为

\[
f ^ {\prime - 1} (f ^ {\prime} (\mathbf {L})) = \mathbf {L}. (\mathbf {G} ^ {\prime} \cap \mathbf {N}),
\]

它定义了……的同构 \(\mathbf{G}' / \mathbf{L}\) . ( \(\mathbf{G}' \cap \mathbf{N}\) ) 到 \(f(\mathbf{G}') / f(\mathbf{L})\) 。根据命题7，(b)， \(f^{-1}(\mathbf{H}') = \mathbf{G}'\mathbf{N}\) ；若 \(f''\) 是 \(\mathbf{G}'\mathbf{N}\) 到 \(\mathbf{H}'\) 它与……一致 \(f\) 在 \(\mathbf{G}'\mathbf{N}\) 的同态，同态 \(\lambda \circ f''\) 于 \(\mathbf{G}'\mathbf{N}\) 到 \(f(\mathbf{G}') / f(\mathbf{L})\) 是满射且核为 \(f^{-1}(f(\mathbf{L})) = \mathbf{L}\mathbf{N}\) ；这证明了 \(\mathbf{L}\mathbf{N}\) 是 \(\mathbf{G}'\mathbf{N}\) 的一个分拆，并且 \(\lambda \circ f''\) 在取商后定义了 \(\mathbf{G}'\mathbf{N} / \mathbf{L}\mathbf{N}\) 到 \(f(\mathbf{G}') / f(\mathbf{L})\) .

推论3。设 \(f: \mathbf{G} \to \mathbf{H}\) 设为一个带算子的群同态，N为其核，X为G的一个子集，使得 \(f(\mathbf{X})\) 生成 H，且 Y 是 N 的子集并生成 N。则 \(\mathbf{X} \cup \mathbf{Y}\) 生成 N。

设 \(G'\) 令G的稳定子群由以下元素生成 \(X \cup Y\) 。如 \(Y \subset G', N \subset G'\) 。如 \(f(X) \subset f(G'), f(G') = H\) 的自同态，因此 \(G' = f^{-1}(H) = G\) .

注：在命题7的记号下，H的子群的原像是G的子群这一事实，可由以下更一般的结论推出。

上定义了一个处处有定义的合成律。若 \(\mathbf{A}\) 和 \(\mathbf{B}\) 是……的子集 \(\mathbf{H}\) 和 \(f\) 是满射，则

\[
f^{-1} (\mathbf {A}. \mathbf {B}) = f^{-1} (\mathbf {A}). f^{-1} (\mathbf {B}), \quad f^{-1} (\mathbf {A} ^ {- 1}) = f^{-1} (\mathbf {A}) ^ {- 1}.
\]

显然 \(f^{-1}(\mathbf{A})\) . \(f^{-1}(\mathbf{B}) \subset f^{-1}(\mathbf{A}.\mathbf{B})\) 另一方面，如果 \(z \in f^{-1}(\mathbf{A}.\mathbf{B})\) ，存在 \(a \in \mathbf{A}\) 和 \(b \in \mathbf{B}\) 的集合，使得 \(f(z) = ab\) ；由于 \(f\) 是满射，存在 \(x \in \mathbf{G}\) 的集合，使得 \(f(x) = a\) ；记 \(y = x^{-1}z\) , \(f(y) = a^{-1}f(z) = b\) 和 \(z = xy\) 的自同态，因此 \(z \in f^{-1}(\mathbf{A})\) . \(f^{-1}(\mathbf{B})\) 。关系 \(x \in f^{-1}(\mathbf{A}^{-1})\) 等价于 \(f(x) \in \mathbf{A}^{-1}\) ，因此到 \(f(x^{-1}) \in \mathbf{A}\) ，即到 \(x^{-1} \in f^{-1}(\mathbf{A})\) 最后到 \(x \in f^{-1}(\mathbf{A})^{-1}\) .

命题8. 设G是一个带算子的群，且A和B是G的两个稳定子群。假设关系式 \(a \in \mathbf{A}\) 和 \(b \in \mathbf{B}\) 蕴含 \(aba^{-1} \in \mathbf{B}^*\) （换言之，A正规化B）。则AB = BA是G的一个稳定子群，A ∩ B是A的一个正规稳定子群，且B是AB的一个正规稳定子群。A到AB的典范嵌入在取商后定义了A/(A ∩ B)到AB/B的一个同构。

公式

\[
\begin{array}{r l} (a b) (a ^ {\prime} b ^ {\prime}) & = a a ^ {\prime} (a ^ {\prime - 1} b a ^ {\prime}. b ^ {\prime}) \\ (a b) ^ {- 1} & = a ^ {- 1} (a b ^ {- 1} a ^ {- 1}) \\ (a b) ^ {\alpha} & = a ^ {\alpha} b ^ {\alpha} \end{array}
\]

对于 \(a, a' \in \mathbf{A}, b, b' \in \mathbf{B}\) 以及每个算子 \(\alpha\) 在 \(\mathbf{G}\) ，证明 AB 是 \(\mathbf{G}\) 的一个映射。设 \(a \in \mathbf{A}\) 和 \(x \in \mathbf{A} \cap \mathbf{B}\) ；于是 \(axa^{-1} \in \mathbf{B}\) 的稳定子群，由对 \(\mathbf{A}\) 和 \(\mathbf{B}\) ，并且显然 \(axa^{-1}\) 属于 \(\mathbf{A}\) ，因此 \(\mathbf{A} \cap \mathbf{B}\) 在 \(\mathbf{A}\) 的一个映射。设 \(a \in \mathbf{A}\) 和 \(b, b'\) 的序列 \(\mathbf{B}\) 所作的假设；公式 \((ab)b'(ab)^{-1} = a(bb'b^{-1})a^{-1}\) 表明 \(\mathbf{B}\) 在 AB 中是正规的。设 \(\phi\) 为到 \(\mathbf{A}\) AB到AB/B的典范同态的 \(\phi(a) = a\mathbf{B}\) 因此核为 \(\phi\) 等于 \(\mathbf{A} \cap \mathbf{B}\) 一般不一定是环同态。显然 \(\phi\) 是满射，因此定义了……的同构 \(\mathbf{A}/(\mathbf{A} \cap \mathbf{B})\) 到AB/B上。

定理4. 设G是一个带算子的群，且N是G的一个正规稳定子群。

\begin{enumerate}
\def\labelenumi{(\alph{enumi})}
\item
  映射 \(\mathbf{G}' \mapsto \mathbf{G}' / \mathbf{N}\) 是稳定子群集合的一个双射 \(\mathbf{G}\) ，包含 \(\mathbf{N}\) 到稳定子群的集合上 \(\mathbf{G} / \mathbf{N}\) .
\item
  设 \(\mathbf{G}'\) 是 \(\mathbf{G}\) ，包含 \(\mathbf{N}\) 。对于 \(\mathbf{G}' / \mathbf{N}\) 在……中正规 \(\mathbf{G} / \mathbf{N}\) ，必要且充分的条件是 \(\mathbf{G}'\) 在……中正规 \(\mathbf{G}\) 以及这些群 \(\mathbf{G} / \mathbf{G}'\) 和 \((\mathbf{G} / \mathbf{N}) / (\mathbf{G}' / \mathbf{N})\) 于是它们是同构的。
\item
  设 \(\mathbf{G}'\) 是 \(\mathbf{G}\) 。则 \(\mathbf{G}'\mathbf{N}\) 是 \(\mathbf{G}\) 和 \(\mathbf{N}\) 在 \(\mathbf{G}'\mathbf{N}\) 。进一步 \(\mathbf{G}' \cap \mathbf{N}\) 在 \(\mathbf{G}'\) 以及这些群 \(\mathbf{G}' / (\mathbf{G}' \cap \mathbf{N})\) 和 \(\mathbf{G}'\mathbf{N} / \mathbf{N}\) 是同构的。
\end{enumerate}

设 \(f\) 表示 \(G\) 到 \(G/N\) 的分拆集合的双射。对于所有 \(x \in G\) , \(f(x) \in xN\) 因此， \(f(G') = G'/N\) 对于每个子群 \(G'\) 于 \(G\) ，包含 \(N\) 。如 \(f\) 是满射，断言(a)由命题7的推论1得出；类似地，对于等价性“ \(G'\) 正规'' \(\Leftrightarrow\) `` \(G'/N\) 正规''。假设 \(G'\) 是 \(G\) ，包含 \(N\) 由第4条，将命题5应用于 \(Id_G\) ，存在一个同态 \(u\) 于 \(G/N\) 到 \(G/G'\) 上的作用，定义为 \(u(xN) = xG'\) 为所有 \(x \in G\) 。显然 \(u\) 是满射且核为 \(G'/N\) 由此得到所需的同构 \((G/N)/(G'/N)\) 到 \(G/G'\) 。最后，(c)直接由命题8得出。

\subsection{约当-赫尔德定理}

定义9. 带算子群G的一个合成列是有限序列 \((\mathbf{G}_i)_{0\leqslant i\leqslant n}\) G的稳定子群，其中 \(\mathbf{G}_0 = \mathbf{G}\) 和 \(\mathbf{G}_n = \{e\}\) 且使得 \(\mathbf{G}_{i + 1}\) 是由 \(\mathbf{G}_i\) 对于 \(0\leqslant i\leqslant n - 1\) 。这些商 \(\mathrm{G_i / G_{i + 1}}\) 被称为该序列的商。一个合成序列 \(\Sigma^{\prime}\) 被称为比一个合成序列更细 \(\Sigma\) 中正规，若 \(\Sigma\) 是从中取出的一个序列 \(\Sigma^{\prime}\) .

上定义了一个处处有定义的合成律。若 \((\mathbf{G}_i)_{0\leqslant i\leqslant n}\) 和 \((\mathbf{H}_j)_{0\leqslant j\leqslant m}\) 分别是两个带算子群的合成群列 \(\mathbf{G}\) 和 \(\mathbf{H}\) ，则称它们是等价的，如果 \(m = n\) 且存在一个排列 \(\phi\) 区间的 \([0,n - 1]\) 于 \(\mathbf{N}\) ，使得带算子的群 \(\mathbf{G}_i / \mathbf{G}_{i + 1}\) 和 \(\mathrm{H}_{\phi (i)} / \mathrm{H}_{\phi (i) + 1}\) 对所有 \(i\) .

注意，一般来说，取自合成列的序列 \((\mathbf{G}_{i})\) 不是合成列，因为对于 \(j > i + 1\) , \(G_{j}\) 一般不构成 \(G_{i}\) .

定理5（施赖埃尔）。给定两个合成列 \(\Sigma_{1}, \Sigma_{2}\) 带算子群G的两个合成列，存在两个等价的合成列 \(\Sigma_{1}^{\prime}, \Sigma_{2}^{\prime}\) ，它们分别加细于 \(\Sigma_{1}\) 和 \(\Sigma_{2}\) .

设 \(\Sigma_{1} = (\mathrm{H}_{i})_{0\leqslant i\leqslant n}\) 和 \(\Sigma_{2} = (\mathrm{K}_{j})_{0\leqslant j\leqslant p}\) 设为两个给定的合成列，分别具有 \(n + 1\) 和 \(p + 1\) 项；我们将看到合成列 \(\Sigma_1^\prime\) 可以由插入形成 \(p - 1\) 子群 \(\mathbf{H}_{i,j}^{\prime}\) \((1 \leqslant j \leqslant p - 1)\) 介于 \(H_{i}\) 和 \(H_{i+1}\) 对于 \(0 \leqslant i \leqslant n - 1\) 与该序列之间 \(\Sigma_{2}'\) 通过插入 \(n - 1\) 子群 \(K_{j,i}'\) ( \(1 \leqslant i \leqslant n - 1\) ) 于 \(K_{j}\) 和 \(K_{j+1}\) 对于 \(0 \leqslant j \leqslant p - 1\) 之间；从而将得到 G 的两列 \(pn + 1\) 稳定子群；通过适当选择所插入的稳定子群，我们将证明这些列是等价的合成列。

为此，注意 \(\mathbf{H}_i \cap \mathbf{K}_j\) 是 \(\mathbf{H}_i\) 以及 \(\mathbf{K}_j\) ，因此（定理4） \(\mathbf{H}_{i+1} \cdot (\mathbf{H}_i \cap \mathbf{K}_j)\) 是 \(\mathbf{H}_i\) ，包含 \(\mathbf{H}_{i+1}\) 和 \(\mathbf{K}_{j+1} \cdot (\mathbf{H}_i \cap \mathbf{K}_j)\) 是 \(\mathbf{K}_j\) ，包含 \(\mathbf{K}_{j+1}\) 。如果我们写 \(\mathbf{H}_{i,j}' = \mathbf{H}_{i+1} \cdot (\mathbf{H}_i \cap \mathbf{K}_j)\) 和 \(\mathbf{K}_{j,i}' = \mathbf{K}_{j+1} \cdot (\mathbf{H}_i \cap \mathbf{K}_j)\) , \(\mathbf{H}_{i,j+1}'\) 是 \(\mathbf{H}_{i,j}'\) ( \(0 \leqslant j \leqslant p-1\) ) 且 \(\mathbf{K}_{j,i+1}'\) 是 \(\mathbf{K}_{j,i}'\) ( \(0 \leqslant i \leqslant n-1\) ）。此外 \(\mathbf{H}_{i,0}' = \mathbf{H}_i\) , \(\mathbf{H}_{i,p}' = \mathbf{H}_{i+1}\) , \(\mathbf{K}_{j,0}' = \mathbf{K}_j\) 和 \(\mathbf{K}_{j,p}' = \mathbf{K}_{j+1}\) 为了证明该定理，只需证明 \(\mathbf{H}_{i,j+1}'\) （相应地 \(\mathbf{K}_{j,i+1}'\) ) 是 \(\mathbf{H}_{i,j}'\) （相应地 \(\mathbf{K}_{j,i}'\) ) 的一个正规稳定子群，并且商群 \(\mathbf{H}_{i,j}'/\mathbf{H}_{i,j+1}'\) 和 \(\mathbf{K}_{j,i}'/\mathbf{K}_{j,i+1}'\) 是同构的（ \(0 \leqslant i \leqslant n-1\) , \(0 \leqslant j \leqslant p-1\) ）。这由以下引理取 \(\mathbf{H} = \mathbf{H}_i\) , \(\mathbf{H}' = \mathbf{H}_{i+1}\) , \(\mathbf{K} = \mathbf{K}_j\) , \(\mathbf{K}' = \mathbf{K}_{j+1}\) .

引理1（Zassenhaus）。设 H 和 K 是带算子群 G 的两个稳定子群，H' 和 K' 分别是 H 和 K 的正规稳定子群；则 H'.(H ∩ K') 是 H'.(H ∩ K) 的一个正规稳定子群，K'.(K ∩ H') 是 K'.(K ∩ H) 的一个正规稳定子群，并且带算子的商群 (H'.(H ∩ K))/(H'.(H ∩ K')) 与 (K'.(K ∩ H))/(K'.(K ∩ H')) 是同构的。

由定理4， \(\mathbf{H}' \cap \mathbf{K} = \mathbf{H}' \cap (\mathbf{H} \cap \mathbf{K})\) 是 \(\mathbf{H} \cap \mathbf{K}\) ；类似地 \(\mathbf{K}' \cap \mathbf{H}\) 是 \(\mathbf{K} \cap \mathbf{H}\) 的一个正规稳定子群；因此（第6节，推论2） \((\mathbf{H}' \cap \mathbf{K})(\mathbf{K}' \cap \mathbf{H})\) 是 \(\mathbf{H} \cap \mathbf{K}\) 。将定理4应用于群 \(\mathbf{H}\) ,

\[
\mathbf {H} ^ {\prime}. (\mathbf {H} ^ {\prime} \cap \mathbf {K}). (\mathbf {K} ^ {\prime} \cap \mathbf {H}) = \mathbf {H} ^ {\prime}. (\mathbf {H} \cap \mathbf {K} ^ {\prime})
\]

是 \(\mathbf{H}'\) . (H ∩ K) 和商群

\[
(\mathbf {H} ^ {\prime}. (\mathbf {H} \cap \mathbf {K})) / (\mathbf {H} ^ {\prime}. (\mathbf {H} \cap \mathbf {K} ^ {\prime}))
\]

（使用复合

\[
(\mathbf {H} \cap \mathbf {K}) / ((\mathbf {H} ^ {\prime} \cap \mathbf {K}). (\mathbf {K} ^ {\prime} \cap \mathbf {H})).
\]

在最后一个商中，H 与 H' 一方面，K 与 K' 另一方面是对称出现的；将它们互换即可得到所述结果。

定义10. 带算子群 G 的一个 Jordan-Hölder 列是一个严格递减的分解列， \(\Sigma\) 使得不存在不同于 \(\Sigma\) 且比其更细的严格递减分解列， \(\Sigma\) .

命题9. 对于 G 的一个分解列，它是 Jordan-Hölder 列的充分必要条件是：该列的所有商都是单的。

一个分解列是严格递减的，当且仅当它的相继商中没有一个是单位元。若一个严格递减的合成列 \(\Sigma\) 不是 Jordan-Hölder 列，则存在一个严格递减的合成列， \(\Sigma'\) 它比 \(\Sigma\) 更细且不同于 \(\Sigma\) 。因此存在两个相继项 \(G_i\) , \(G_{i+1}\) 于 \(\Sigma\) 它们在 \(\Sigma'\) 中不是相继的；设 H 是紧随 \(G_i\) 在 \(\Sigma'\) 之后的第一个项；H 是 \(G_i\) 的一个正规稳定子群，且包含 \(G_{i+1}\) 且与后者不同；因此 H/ \(G_{i+1}\) 是 \(G_i/G_{i+1}\) ，与后者不同且不同于单位元；因此 \(G_i/G_{i+1}\) 不是单群。反之，如果 \(\Sigma\) 是一个严格递减的合成列，其一个商 \(G_i/G_{i+1}\) 不是单的，这个商群包含一个不同于自身的正规稳定子群。 \(\{e\}\) ，其在 \(G_i\) 是 G 的一个正规稳定子群 H， \(G_i\) ，不同于 \(G_i\) 和 \(G_{i+1}\) (定理4)；只需在H之间插入 \(G_i\) 和 \(G_{i+1}\) 为了得到一个严格递减的合成列，且不同于 \(\Sigma\) 且比其更细的严格递减分解列， \(\Sigma\) .

定理6（Jordan-Hölder）：一个带算子群的任意两条Jordan-Hölder序列是等价的。

设 \(\Sigma_{1}, \Sigma_{2}\) 设G是一个带算子群的两条Jordan-Hölder序列；通过应用定理5，两条等价的合成序列 \(\Sigma_{1}^{\prime}, \Sigma_{2}^{\prime}\) 分别比 \(\Sigma_{1}\) 和 \(\Sigma_{2}\) 更细；由于后者是Jordan-Hölder序列， \(\Sigma_{1}^{\prime}\) 与……相同 \(\Sigma_{1}\) 或者是由它通过重复某些项而得到的；其商的序列 \(\Sigma_{1}^{\prime}\) 是由……的推导而来 \(\Sigma_{1}\) 通过插入若干与群同构的项 \(\{e\}\) ；因为 \(\Sigma_{1}\) 严格递减，因此……的商序列 \(\Sigma_{1}\) 是由 \(\Sigma_{1}^{\prime}\) 的推导而来，通过抑制后者中所有同构于 \(\{e\}\) 类似地，对于 \(\Sigma_{2}\) 和 \(\Sigma_{2}^{\prime}\) 。由于商序列 \(\Sigma_{1}^{\prime}\) 和 \(\Sigma_{2}^{\prime}\) （在同构意义下）仅在各项的顺序上有所不同，这一点同样适用于 \(\Sigma_{1}\) 和 \(\Sigma_{2}\) 定理得证。

推论。设G是一个具有算子且存在若尔当-赫尔德列的群。若 \(\Sigma\) 对于G的任何严格递减合成列，存在一个比其更细的Jordan-Hölder列。 \(\Sigma\) .

设 \(\Sigma_0\) 是 \(G\) ；根据定理5，存在两个等价的合成序列， \(\Sigma'\) 和 \(\Sigma_0'\) 分别比……更细 \(\Sigma\) 和 \(\Sigma_0\) 定理6的论证表明，通过从 \(\Sigma'\) 重复，一个序列 \(\Sigma''\) 得到，这与……等价。 \(\Sigma_0\) 因此它是一个Jordan-Hölder序列，因为其所有商都是单的（命题9）。 \(\Sigma\) 是严格递减的， \(\Sigma''\) 比……更细 \(\Sigma\) ，由此得出推论。

注：带算子的群并不总是具有Jordan-Hölder列；一个例子由加法群给出。 \(\mathbf{Z}\) 有理整数的序列：该序列 \((2^{n}.\mathbf{Z})_{n\geqslant 0}\) 是 \(\mathbf{Z}\) ；对于所有 \(p\) 的（正规）子群的严格递减无穷序列，其前 \(p\) 项与该群一起构成 \(\{0\}\) 一个严格递减的合成列；如果Z存在一个Jordan-Hölder列，那么它至少会有 \(p + 1\) 由定理6的推论，各项；荒谬，因为p是任意的。

另一方面，在每一个带算子的有限群G中都存在一个Jordan-Hölder合成列：若 \(G \neq \{e\}\) 在G的异于G的正规稳定子群中，设 \(H_{1}\) 是极大子群；类似地定义 \(H_{n+1}\) 通过归纳法，作为 \(H_{n}\) 不同于 \(H_{n}\) 的正规子群集合中的极大元素，当 \(H_{n} \neq \{e\}\) ; 这些阶的序列 \(H_{n}\) 严格递减，因此存在一个n使得 \(H_{n} = \{e\}\) 由G以及 \(H_{i} (1 \leqslant i \leqslant n)\) 按其构造，这是一个Jordan-Hölder序列。

定义11. 设G是一个带算子的群；G的长度是使得存在G的严格递减合成列（G \(_{i}\) ) \(_{0 \leq i \leq n}\) .

若G具有Jordan-Hölder列，则G的长度为该列的相继商的数量，这由定理6的推论得出。若G不具有Jordan-Hölder列，则其长度为无穷；根据命题9，对于G的每个严格递减列，都存在一个严格更细的严格递减列。由单位元构成的群是唯一长度为0的带算子群。一个带算子群是单群当且仅当其长度为1。

设G是一个带算子的群，H是G的正规稳定子群，K是商群G/H，且 \(\pi: G \rightarrow K\) 典范同态。设

\[
\Sigma^ {\prime} = \left(\mathrm{H} _ {i}\right) _ {0 \leqslant i \leqslant n}
\]

是 H 的一个分解列，且 \(\Sigma'' = (\mathbf{K}_j)_{0 \leqslant j \leqslant p}\) 设K的一个合成列，记作 \(G_i = \pi^{-1}(\mathbf{K}_i)\) 对于 \(0 \leqslant i \leqslant p\) 和 \(G_i = H_{i-p}\) 对于 \(p \leqslant i \leqslant n + p\) ，一个合成列 \(\Sigma = (G_i)_{0 \leqslant i \leqslant n+p}\) 就得到了。其商序列 \(\Sigma\) 是通过将 \(\Sigma''\) 的商序列与 \(\Sigma'\) 的一个正规子群。若 \(\Sigma'\) 和 \(\Sigma''\) 的商序列并列而得到的， \(\Sigma\) 是G的Jordan-Hölder序列，由命题9可知。若H或K具有任意长度的合成序列，则G亦然。我们已证明：

命题10. 设G是一个带算子的群，H是G的一个正规稳定子群。G的长度等于H与G/H的长度之和。

推论. 设G是一个带算子的群，且 \((\mathrm{G}_{i})_{0\leqslant i\leqslant n}\) G 的一个合成列。G 的长度是 \(G_{i}/G_{i+1}, 0 \leqslant i \leqslant n - 1\) .

若G与G'是同构的带算子群，且G具有Jordan-Hölder列，则G'亦然，且G与G'的Jordan-Hölder列等价。然而，不同构的群也可能具有等价的Jordan-Hölder列；例如Z/4Z与(Z/2Z) × (Z/2Z)即如此，参见第10号。

\subsection{乘积与纤维积}

设 \((\mathbf{G}_i)_{i\in I}\) 是一族带算子的群。设 \(\mathbf{G}\) 为 \(\mathbf{G}_i\) 的乘积幺半群。考虑 \(\Omega\) 在 \(\mathbf{G}\) 上的作用，定义为

\[
((x _ {i}) _ {i \in I}) ^ {\alpha} = (x _ {i} ^ {\alpha}) _ {i \in I} (\alpha \in \Omega , x _ {i} \in \mathrm{G} _ {i}).
\]

在此结构下，G 是一个带算子的群。对所有 \(i \in I\) ，投影映射 \(pr_{i}: G \to G_{i}\) 是带算子群的同态。

定义12. 上述定义的带算子群 \(\mathbf{G} = \prod_{i\in I}\mathbf{G}_i\) 是 \(\mathbf{G}_i\) 的带算子群乘积。映射 \(\mathbf{pr}_i\colon \mathbf{G}\to \mathbf{G}_i\) 称为投影同态。

带算子群乘积的一个特殊情形是群 \(G^{E}\) ，它由集合E到带算子群G的映射组成，其运算定义为：

\[
\begin{array}{r l} (f g) (x) & = f (x) g (x) \quad (f, g \in \mathbf {G} ^ {E}, x \in \mathbf {E}) \\ f ^ {\alpha} (x) & = f (x) ^ {\alpha} \quad (f \in \mathbf {G} ^ {E}, \alpha \in \Omega , x \in \mathbf {E}). \end{array}
\]

设 \((\phi_{i} \colon H \to G_{i})_{i \in I}\) 是带算子群的同态族。映射 \(h \mapsto (\phi_{i}(h))_{i \in I}\) 的 H 到 \(\prod_{i \in I} G_{i}\) 是带算子的群同态。它是唯一的同态。 \(\Phi \colon H \to \prod_{i \in I} G_{i}\) 满足 \(\mathbf{pr}_{i} \circ \Phi = \phi_{i}\) 对所有i成立。这证明了“带算子的积群”这一名称的合理性（《集合论》，IV，§2，第4节）。

设 \((\phi_i: \mathrm{H}_i \to \mathrm{G}_i)_{i \in I}\) 是带算子群的同态族。映射 \(\prod_{i \in I} \phi_i: (h_i)_{i \in I} \mapsto (\phi_i(h_i))_{i \in I}\) 于 \(\prod_{i \in I} \mathrm{H}_i\) 到 \(\prod_{i \in I} \mathrm{G}_i\) 是带算子群的同态。

命题11。设 \((\phi_i: \mathrm{H}_i \to \mathrm{G}_i)_{i \in I}\) 设为一族带算子的群同态，并令 \(\Phi = \prod_{i \in I} \phi_i\) 那么：

\begin{enumerate}
\def\labelenumi{(\alph{enumi})}
\item
  \(\operatorname{Ker}(\Phi) = \prod_{i \in I} \operatorname{Ker}(\phi_i)\) ；特别地，如果所有 \(\phi_i\) 都是单射， \(\Phi\) 是单射。
\item
  \(\operatorname{Im}(\Phi) = \prod_{i\in I}\operatorname{Im}(\phi_i)\) ；特别地，如果所有 \(\phi_i\) 都是满射， \(\Phi\) 是满射。
\end{enumerate}

这是显然的。

特别是，让 \((\mathbf{G}_{i})_{i\in\mathbf{I}}\) 设为一族带算子的群，且对所有 i，令 \(\mathbf{H}_i\) 是……的一个稳定（相应地，正规稳定）子群 \(\mathbf{G}_i\) 。乘积 \(\prod_{i\in I}\mathbf{H}_i\) 是稳定（相应地，正规稳定）子群 \(\prod_{i\in I}\mathbf{G}_i\) 以及 \(\prod_{i\in I}\mathbf{G}_i\) 到 \(\prod_{i\in I}(\mathbf{G}_i / \mathbf{H}_i)\) 定义了在过渡到商时的一个同构

\[
\left(\prod_ {i \in I} G _ {i}\right) / \left(\prod_ {i \in I} H _ {i}\right)
\]

到 \(\prod_{i\in I}(\mathrm{G}_i / \mathrm{H}_i)\) 。例如，设 J 是 I 的一个子集。子群 \(\mathbf{G}_{J}\) 于 \(\prod_{i\in I}\mathbf{G}_i\) 由 \((x_{i})_{i\in I}\) 的集合，使得 \(x_{i} = e_{i}\) 对于 \(i\notin J\) 是一个正规稳定子群。映射 \(\iota_{J}\) 它将 \(x = (x_{j})_{j\in J}\) 元素 \(y = (y_{i})_{i\in I}\) 的集合，使得 \(y_{i} = e_{i}\) 对于 \(i\notin J\) 和 \(y_{i} = x_{i}\) 对于 \(i\in J\) ，是 \(\prod_{j\in J}\mathbf{G}_j\) 到 \(\mathbf{G}_{J}\) 该映射 \(\operatorname{pr}_{I - J}\) 定义了在过渡到商群时的同构 \(\theta_J\) 商群的 \(\mathbf{G} / \mathbf{G}_J\) 到 \(\prod_{i\in I - J}\mathbf{G}_i\) 。复合 \(\operatorname{pr}_J\circ \iota_J\) 是恒等映射 \(\prod_{j\in J}\mathbf{G}_j\) . \(\mathbf{G} / \mathbf{G}_J\) 通常与 \(\prod_{i\in I - J}\mathbf{G}_i\) 等同，因为 \(\theta_J\) 和 \(\prod_{i\in J}\mathbf{G}_i\) 使得 \(\mathbf{G}_J\) 等同，因为 \(\iota_J\) .

上定义了一个处处有定义的合成律。若 \(J_{1}\) 和 \(J_{2}\) 是I的不相交子集，根据定义可知，I的每个元素 \(G_{J_{1}}\) 与……的每个元素交换 \(G_{J_{2}}\) .

定义13。设G是一个带算子的群，且 \((\mathrm{H}_{i})_{i\in I}\) G的一族正规稳定子群。设 \(p_{i}\colon G\to G/H_{i}\) 是典范同态。称G为商群族的内部积（或积） \((\mathrm{G}/\mathrm{H}_{i})\) 如果同态 \(g\mapsto(p_{i}(g))_{i\in I}\) 是 G 到 \(\prod_{i\in I}G/H_{i}\) .

设 G 和 H 是带算子的群，并设 \(\phi\) 和 \(\psi\) 是 G 到 H 的两个同态。满足 \(x\) 的 G 中的元素集合 \(\phi(x) = \psi(x)\) 是 G 的一个稳定子群，称为 \(\phi\) 和 \(\psi\) 的重合群。特别地，设 \(\phi_1: G_1 \to H\) 和 \(\phi_2: G_2 \to H\) 是带算子群的同态；这些同态 \(\phi_1 \circ \mathrm{pr}_1\) 和 \(\phi_2 \circ \mathrm{pr}_2\) 于 \(G_1 \times G_2\) 到 H 中的重合群称为 \(G_1\) 和 \(G_2\) 在 H 上相对于 \(\phi_1\) 和 \(\phi_2\) 的同态。它被记为 \(G_1 \times_H G_2\) 当没有歧义时 \(\phi_1\) 和 \(\phi_2\) 以及限制条件 \(p_1\) 和 \(p_2\) 于 \(\mathrm{pr}_1\) 和 \(\mathrm{pr}_2\) 到 \(G_1 \times_H G_2\) 也称为投影同态。 \(\phi_1 \circ p_1 = \phi_2 \circ p_2\) 。 \(G_1 \times_H G_2\) 是有序对 \((g_1, g_2) \in G_1 \times G_2\) 的集合，使得 \(\phi_1(g_1) = \phi_2(g_2)\) 的一个正规子群。若 \(f_i\) 是带算子群 K 到 \(G_i (i = 1, 2)\) 和 \(\phi_1 \circ f_1 = \phi_2 \circ f_2\) 的同态，则存在且仅存在一个同态 \(f\) 的K到 \(G_1 \times_H G_2\) 的集合，使得 \(f_i = p_i \circ f\) 对于 \(i = 1, 2\) .

\subsection{受限和}

设 \((\mathbf{G}_i)_{i\in I}\) 是带算子群的一个族，且对于 \(i\in I\) ，设 \(\mathbf{H}_i\) 是 \(\mathbf{G}_i\) 的子集，使得 \(\prod_{i\in I}\mathbf{G}_i\) 由 \((x_i)_{i\in I}\) 的集合 \(i\in I\) 使得 \(x_i\notin H_i\) 是有限的是一个稳定子群 \(\prod_{i\in I}\mathbf{G}_i\) 等于 \(\prod_{i\in I}\mathbf{G}_i\) 中正规，若 \(I\) 是有限的。它被称为受限和，即 \(\mathbf{G}_i\) 关于 \(\mathbf{H}_i\) 。当对所有 \(i\) 除有限个外 \(\mathbf{H}_i\) 是 \(\mathbf{G}_i\) ，受限和是乘积的正规稳定子群。当对所有 \(i\) ，子群 \(\mathbf{H}_i\) 约化为 \(\mathbf{G}_i\) 的单位元时， \(\mathbf{G}_i\) 关于 \(\mathbf{H}_i\) 的直和简称为 \(\mathbf{G}_i\) 的受限和，有时记为 \(\prod_{i\in I}\mathbf{G}_i\) 的分拆集合的双射。对于所有 \(i_0\in I\) ，映射 \(\iota_{i_0}:\mathrm{G}_{i_0}\to \prod_{i\in I}\mathrm{G}_i\) 上的作用，定义为 \(\iota_{i_0}(x) = (x_i)_{i\in I}\) 的x和y，其中 \(x_{i_0} = x\) 和 \(x_i = e_i\) 中正规，若 \(i\neq i_0\) ，是带算子群的单射同态，称为典范嵌入。 \(\mathbf{G}_i\) 与稳定子群 \(\operatorname {Im}(\iota_i)\) 等同。子群 \(\mathbf{G}_i\) 是正规的。对于 \(i\neq j\) ，的元素 \(\mathbf{G}_i\) 和 \(\mathbf{G}_j\) 交换且 \(\mathrm{G}_i\cap \mathrm{G}_j = \{e\}\) 。群 \(\prod_{i\in I}\mathrm{G}_i\) 由集合生成 \(\bigcup_{i\in I}\mathrm{G}_i\) .

命题12. 设 \((\phi_i: \mathbf{G}_i \to \mathbf{K})_{i \in I}\) 是一族带算子的群同态，使得对所有 \(i \in I\) 和 \(j \in I\) 使得 \(i \neq j\) , \(x \in \mathbf{G}_i\) , \(y \in \mathbf{G}_j\) ，元素 \(\phi_i(x)\) 和 \(\phi_j(y)\) 于 \(\mathbf{K}\) 交换；存在且仅存在一个带算子的群同态 \(\Phi\) 于 \(\prod_{i \in I} \mathbf{G}_i\) 到 \(\mathbf{K}\) 的集合，使得 \(\phi_i = \Phi \circ \iota_i\) 为所有 \(i \in I\) 。对于每个元素 \(x = (x_i)_{i \in I}\) 于 \(\prod_{i \in I} \mathbf{G}_i\) , \(\Phi(x) = \prod_{i \in I} \phi_i(x_i)\) .

上定义了一个处处有定义的合成律。若 \(\Phi\) 和 \(\Phi'\) 是问题的解，它们在 \(\bigcup_{i\in I}G_i\) 上重合，从而在 \(\prod_{i\in I}G_i\) 上重合，由此得到 \(\Phi\) 的唯一性。我们现在证明 \(\Phi\) 的存在性：对每个元素 \(x = (x_i)_{i\in I}\) 于 \(\prod_{i\in I}G_i\) ，设 \(\Phi(x) = \prod_{i\in I}\phi_i(x_i)\) ( \(\S 1\) , 第5号）。显然 \(\Phi \circ \iota_i = \phi_i\) 为所有 \(i\) 和 \(\Phi\) 与位似交换；公式 \(\Phi(xy) = \Phi(x)\Phi(y)\) 由 \(\S 1\) ，第5号，公式(9)。

定义14. 设G是一个带算子的群，且 \((\mathbf{H}_i)_{i\in I}\) 是G的一族稳定子群。若G的每个元素 \((\mathbf{H}_i)\) 且从 \(\mathbf{H}_i\) 与以下每个元素可交换 \(\mathbf{H}_j\) 对于 \(j\neq i\) 到G的唯一同态，其在每个 \(\prod_{i\in I}\mathbf{H}_i\) 上的限制为典范嵌入，则该同态为同构，则称G为子群族 \(\mathbf{H}_i\) 的内受限和（或受限和）。

当I有限时，我们滥用语言，也称内直积（或直积，或积）代替内受限和。G的每个稳定子群H，若存在一个稳定子群 \(H'\) 的G，使得G是H与 \(H'\) 的直积，称为G的一个直因子。

命题13。设G是一个带算子的群，且 \((\mathbf{H}_i)_{i\in I}\) G的一族稳定子群，使得G的每个元素 \(\mathbf{H}_i\) 与以下每个元素可交换 \(\mathbf{H}_j\) 对于 \(j\neq i\) 要使 G 成为子群族 \((\mathbf{H}_i)_{i\in I}\) 的受限和，必要且充分条件是 G 的每个元素 \(x\) 都能唯一地表示为 \(\prod_{i\in I}y_i\) 的x和y，其中 \((y_i)_{i\in I}\) 是 G 中元素的一个有限支撑族，且 \(y_i\in \mathbf{H}_i\) 为所有 \(i\) .

显然。

命题14。设G是一个带算子的群，且 \((\mathbf{H}_i)_{i\in I}\) G 的一个有限稳定子群族。要使 G 成为该子群族的限制和 \((\mathbf{H}_i)\) ，必要且充分的条件是每个 \(\mathbf{H}_i\) 是正规的，且G是商群的乘积 \((\mathbf{G} / \mathbf{H}^i)\) 的x和y，其中 \(\mathbf{H}^i\) 是由 \(\mathbf{H}_j\) 对于 \(j\neq i\) .

该条件是显然必要的。反之，假设G是……的乘积。 \(K_{i} = G/H^{i}\) 并将 G 与 \(K_{i}\) 。则 \(H_{i}\) 被等同于 \(K_{i}\) 的一个子群，因此对于 \(i \neq j\) ，中的每个元素 \(H_{i}\) 与以下每个元素可交换 \(H_{j}\) ；另一方面， \(H^{i}\) 被等同于 \(K_{j}\) 对于 \(j \neq i\) ，因此 \(H_{i} = K_{i}\) 对所有 i 成立，且 G 是 \(H_{i}\) .

命题15。设G是一个带算子的群，且 \((\mathbf{H}_i)_{1\leqslant i\leqslant n}\) G 的一个正规稳定子群序列，使得

\[
\left(\mathrm{H} _ {1} \mathrm{H} _ {2} \dots \mathrm{H} _ {i}\right) \cap \mathrm{H} _ {i + 1} = \{e \} \text { for } 1 \leqslant i \leqslant n - 1,
\]

表示集合 \(\mathbf{H}_1\mathbf{H}_2\ldots \mathbf{H}_n\) 是 \(\mathbf{G}\) ，即 \(\mathbf{H}_i\) .

通过对 \(n\) 进行归纳，这立即归结为对 \(n = 2\) 证明该命题。我们首先证明，若 \(x \in \mathrm{H}_1\) 和 \(y \in \mathrm{H}_2\) , \(x\) 和 \(y\) 是可交换的；因为 \(xyx^{-1}y^{-1} = (xyx^{-1})y^{-1} = x(yx^{-1}y^{-1})\) 且因此（由于 \(\mathrm{H}_1\) 和 \(\mathrm{H}_2\) 是正规的） \(xyx^{-1}y^{-1} \in \mathrm{H}_1 \cap \mathrm{H}_2\) ，即 \(xyx^{-1}y^{-1} = e\) ，由假设可得。此外 \(\mathrm{H}_1\mathrm{H}_2\) 是 \(G\) 的子集，且它在 \(G\) 的位似变换下稳定。由此（根据第3节命题1）可得 \(\mathrm{H}_1\mathrm{H}_2\) 是 \(G\) ，并且容易验证该子群是正规的。最后假设 \(xy = x'y'\) ，其中 \(x \in \mathrm{H}_1\) , \(x' \in \mathrm{H}_1\) , \(y \in \mathrm{H}_2\) , \(y' \in \mathrm{H}_2\) ；于是 \(x'^{-1}x = y'y^{-1}\) ，因此 \(x'^{-1} \in \mathrm{H}_1 \cap \mathrm{H}_2 = \{e\}\) , \(x' = x\) 并且类似地 \(y' = y\) ; \(\mathrm{H}_1\mathrm{H}_2\) 因此是 \(\mathrm{H}_1\) 和 \(\mathrm{H}_2\) .

当所考虑的群是交换群时，使用术语“直和”而非“直积”。

\subsection{单生成群}

设 \(a \in Z\) ；因为 \(aZ\) 是 \(Z\) ，元素之间的关系 \(x, y\) 于 \(Z\) 它表明“存在 \(z \in Z\) 的集合，使得 \(x - y = az\) ”是一个等价关系，我们一致同意，一劳永逸地将其写为 \(x \equiv y \pmod{a}\) 或 \(x \equiv y(a)\) 并称之为模 \(a\) 的同余。将 \(a\) 中条件的扩张 \(-a\) 由此得到一个等价关系，因此可以假设 \(a \geqslant 0\) ；对于 \(a = 0\) , \(x \equiv y(0)\) 意味着 \(x = y\) ，因此只有当 \(a \neq 0\) 时才得到与相等不同的关系：因此在下文中我们将假设 \(a > 0\) 除非另有说明。

对于 \(a > 0\) ，商 \(\mathbf{Z}\) 由同余关系 \(x \equiv y(a)\) ，即群 \(\mathbf{Z}/a\mathbf{Z}\) ，被称为有理整数模 \(a\) .

命题16. 设 \(a\) 为整数 \(>0\) . 整数 \(r\) 的集合，使得 \(0 \leqslant r < a\) 构成该等价关系的一个代表系 \(x \equiv y \pmod{a}\) 在 \(\mathbf{Z}\) .

上定义了一个处处有定义的合成律。若 \(x\) 是一个整数 \(\geqslant 0\) 存在（集合论，III，§5，第6号）整数 \(q\) 和 \(r\) 的集合，使得 \(x = aq + r\) 和 \(0 \leqslant r < a\) 和 \(x \equiv r \pmod{a}\) 的一个正规子群。若 \(x\) 是一个整数 \(\leqslant 0\) ，整数 \(-x\) 是 \(\geqslant 0\) 且由上述存在一个整数 \(r\) 的集合，使得 \(0 \leqslant r < a\) 和 \(-x \equiv r \pmod{a}\) 。写 \(r' = 0\) 中正规，若 \(r = 0\) 和 \(r' = a - r\) 中正规，若 \(r > 0\) ，则

\[
x \equiv - r \equiv r ^ {\prime} \pmod {a}
\]

和 \(0 \leqslant r' < a\) 。我们现在证明，如果 \(0 \leqslant r < r' < a\) ，则 \(r \not\equiv r' \pmod{a}\) 的一个子群。现在 \(r' - r < na\) 对于 \(n \geqslant 1\) 和 \(r' - r > na\) 对于 \(n \leqslant 0\) 的自同态，因此 \(r' - r \notin a\mathbf{Z}\) .

推论。设 \(a\) 为整数 \(> 0\) 。群 \(\mathbf{Z} / a\mathbf{Z}\) 有理整数模 \(a\) 是一个阶为 \(a\) .

命题17. 设 H 是 Z 的一个子群。存在唯一一个整数 \(a \geqslant 0\) 使得 H = aZ。

上定义了一个处处有定义的合成律。若 \(\mathbf{H} = \{0\}\) ，则 \(\mathbf{H} = 0\mathbf{Z}\) 的两个元素。假设 \(\mathbf{H} \neq \{0\}\) 子群 \(\mathbf{H}\) 有一个元素 \(x \neq 0\) 。则 \(x > 0\) 或 \(-x > 0\) ，因此 \(\mathbf{H}\) 有元素 \(>0\) 的一个映射。设 \(a\) 是最小元素 \(>0\) 在 \(\mathbf{H}\) 子群 \(a\mathbf{Z}\) 生成的 \(a\) 包含于 \(\mathbf{H}\) ；我们证明 \(\mathbf{H} \subset a\mathbf{Z}\) 的一个映射。设 \(y \in \mathbf{H}\) 。根据命题16，存在一个整数 \(r\) 的集合，使得 \(y \equiv r (\text{mod. } a)\) 和 \(0 \leqslant r < a\) 更不用说 \(y \equiv r (\text{mod. } H)\) 的自同态，因此 \(r \in \mathbf{H}\) 但这仅在以下情况下才可能 \(r = 0\) ，从而 \(y \in a\mathbf{Z}\) 的一个分解。整数 \(a\) 是唯一的：如果 \(\mathbf{H} = \{0\}\) ，那么必然 \(a = 0\) ，且若 \(\mathbf{H} \neq \{0\}\) ，整数 \(a\) 是……的阶 \(\mathbf{Z}/\mathbf{H}\) .

定义15. 如果一个群允许由单个元素组成的生成元系，则称其为单生成群。有限单生成群称为循环群。

每个单生成群都是交换群（第3节，命题2的推论2）。单生成群的每个商群都是单生成群（第3节，命题2的推论3）。

加法群 Z 是单生成的：它由 \(\{1\}\) 对于每个正整数 a，群 Z/aZ 是单生成的，因为它是 Z 的商群。

命题18. 一个有限阶的单一生成群 \(a\) （使用复合 \(\mathbf{Z}/a\mathbf{Z}\) 。一个无限单生成群同构于 \(\mathbf{Z}\) .

设 G 是一个单生成群（乘法记号），x 是 G 的一个生成元。恒等元 \(x^{m}x^{n} = x^{m+n}\) (§ 1，第 3 号，公式 (1))表明该映射 \(n \mapsto x^{n}\) 是从Z到G的同态。其像是G的一个包含x的子群，因此就是G。根据第5条定理3，群G同构于Z关于某个子群的商群，该子群必然具有aZ的形式（命题17）。若a \textgreater{} 0时，群G是阶为a的有限群；若a = 0，则群G同构于Z。

命题19。设 \(a\) 为整数 \(>0\) 的一个映射。设 \(\mathbf{H}\) 的一个子群，且包含于 \(\mathbf{Z}/a\mathbf{Z}\) , \(b\) 的顺序 \(\mathbf{H}\) 和 \(c\) 其在……中的索引 \(\mathbf{Z}/a\mathbf{Z}\) 。则 \(a = bc\) , \(\mathbf{H} = c\mathbf{Z}/a\mathbf{Z}\) 和 \(\mathbf{H}\) （使用复合 \(\mathbf{Z}/b\mathbf{Z}\) .

反之，设 \(b\) 和 \(c\) 是两个整数 \(>0\) 的集合，使得 \(a = bc\) 。则 \(a\mathbf{Z} \subset c\mathbf{Z}\) 和 \(c\mathbf{Z}/a\mathbf{Z}\) 是 \(\mathbf{Z}/a\mathbf{Z}\) ，其阶为 \(b\) 且指数为 \(c\) .

由第4号命题的推论4，a = bc。根据第7号定理4，H具有形式 \(H'/aZ\) 的x和y，其中 \(H'\) 是Z的一个子群，且 \(Z/H'\) （使用复合 \((\mathbf{Z}/a\mathbf{Z})/H\) 因此阶为 c。由命题 17 及命题 16 的推论， \(H' = cZ\) 因此 H 是单生成的。最后，由命题 18，H 同构于 Z/bZ。反之，若 a = bc，则 \(aZ \subset cZ\) 对于 \(a \in cZ\) ：商群 \((\mathbf{Z}/a\mathbf{Z})/(c\mathbf{Z}/a\mathbf{Z})\) 同构于 Z/cZ（第 7 节，定理 4），因此阶为 c（第 4 节，命题 4 的推论），指数为 b（第 4 节，命题 4 的推论）。

推论。单生成群的每个子群都是单生成的。

设 \(a\) 和 \(b\) 是两个整数 \(\neq 0\) 。关系 \(b \in a\mathbf{Z}\) 也写作： \(b\) 是 \(a\) 的倍数，并且也 \(a\) 整除 \(b\) 或 \(a\) 是……的因子 \(b\) .

定义16. 一个整数 \(p > 0\) 被称为素数，如果 \(p \neq 1\) 且它没有因子 \(> 1\) 除了 \(p\) .

命题20。一个整数 \(p > 0\) 当且仅当该群 \(\mathbf{Z} / p\mathbf{Z}\) 是单群。

这由命题19得出。

推论。每个交换单群都是素数阶循环群。

设 G 是这样的一个群。则 \(G \neq \{e\}\) ；令 \(a \neq e\) 是G的一个元素。由a生成的子群是正规的，因为G是交换的，它并不缩减为 \(\{e\}\) 因此等于G。所以G是单生成的，从而同构于形如Z/pZ的群，其中p \textgreater{} 0，因为Z不是单群，且p必然是素数。

注记。阶为素数的有限群G必然是循环群。G除G本身和 \(\{e\}\) 外没有其他子群，因此它由每个非单位元素生成。 \(\neq e\) .

引理2。设a为整数 \textgreater0。通过将每个合成列与 \((\mathbf{H}_{i})_{0\leqslant i\leqslant n}\) 在群 Z/aZ 中，序列 \((s_{i})_{1\leqslant i\leqslant n}\) 的x和y，其中 \(s_{i}\) 是……的阶 \(H_{i-1}/H_{i}\) ，从而在 Z/aZ 的合成列与有限序列 \((s_{i})\) 整数 \textgreater0，使得 \(a=s_{1}\ldots s_{n}\) 之间建立一一对应。合成列 \((\mathbf{H}_{i})_{0\leqslant i\leqslant n}\) 是 Jordan-Hölder 列当且仅当 \(s_{i}\) 为素数。

上定义了一个处处有定义的合成律。若 \((\mathbf{H}_i)_{0\leqslant i\leqslant n}\) 是 \(\mathbf{Z} / a\mathbf{Z}\) ，由此，通过对 \(n\) 进行归纳，根据第4号，命题4，可知 \(a = \prod_{i=1}^{n} (\mathbf{H}_{i-1}:\mathbf{H}_i)\) .

反之，设 \((s_i)_{1 \leqslant i \leqslant n}\) 是一个整数序列 \(>0\) 的集合，使得 \(a = s_1 \ldots s_n\) 的一个正规子群。若 \((\mathbf{H}_i)_{0 \leqslant i \leqslant n}\) 是 \(\mathbf{Z} / a\mathbf{Z}\) 的集合，使得 \((\mathbf{H}_{i-1} : \mathbf{H}_i) = s_i\) 对于 \(1 \leqslant i \leqslant n\) ，那么必然 \(((\mathbf{Z} / a\mathbf{Z}) : \mathbf{H}_i) = \prod_{1 \leqslant j \leqslant i} s_j\) 正如通过对 \(i\) 的自同态，因此 \(\mathbf{H}_i = \left( \prod_{j=1}^{i} s_j \right) \mathbf{Z} / a\mathbf{Z}\) 进行归纳所见（命题19），这显示了所讨论映射的单射性。我们现在证明其满射性。写出 \(H_{i}=\left(\prod_{j=1}^{i}s_{j}\right)\mathbf{Z}/a\mathbf{Z}\) 对于 \(0\leqslant i\leqslant n\) ，得到Z/aZ的一个合成列，使得 \((\mathrm{H}_{i-1}:\mathrm{H}_{i})=s_{i}\) （命题19）。第二个断言由命题20及第7号，命题9得出。

设 \(\mathfrak{P}\) 表示素数的集合。

定理7（分解为素因数）。设 \(a\) 设为一个严格正整数。存在且仅存在一个族 \((v_{p}(a))_{p \in \mathfrak{P}}\) 整数 \(\geqslant 0\) 的集合 \(p \in \mathfrak{P}\) 使得 \(v_{p}(a) \neq 0\) 是有限的且

\[
a = \prod_ {p \in \mathfrak {P}} p ^ {v _ {p} (a)}.\tag{5}
\]

当群 \(\mathbf{Z}/a\mathbf{Z}\) 是有限的，则它有一个若尔当-赫尔德序列。引理2随后蕴含 \(a\) 是素数整数的乘积，因此该族的存在性 \((v_{p}(a))\) ；进一步地，对于每一个族 \((v_{p}(a))_{p\in\mathfrak{P}}\) 满足定理7的条件，整数 \(v_{p}(a)\) 是，对于所有 \(p\in\mathfrak{P}\) ，等于一个Jordan-Hölder序列的因子数 \(\mathbf{Z}/a\mathbf{Z}\) 同构于 \(\mathbf{Z}/p\mathbf{Z}\) （引理2）。该族的唯一性 \((v_{p}(a))\) 因此由Jordan-Hölder定理（第7号，定理6）得出。

推论。设 \(a\) 和 \(b\) 是两个整数 \(>0\) 。则 \(v_{p}(ab) = v_{p}(a) + v_{p}(b)\) 。对于 \(a\) 整除 \(b\) ，必要且充分的条件是 \(v_{p}(a) \leqslant v_{p}(b)\) 对于每个素数 \(p\) .

在任意群 G 中，如果由元素 \(x \in G\) 生成的（单生成）子群的阶为有限数 d，则称 x 为 d 阶元素；因此数 d 是满足 \textgreater0，使得 \(x^{d} = e\) 的最小整数；如果由 x 生成的子群是无限的，则称 x 为无限阶元素。这些定义连同命题 4（第 4 节）特别蕴含：在有限群 G 中，G 的每个元素的阶都是 G 的阶的因子。

命题21。设G是一个阶为n的有限群， \(x^{n} = e\) 为所有 \(x \in G\) .

上定义了一个处处有定义的合成律。若 \(p\) 是……的阶 \(x\) ，则 \(n = pq\) ，其中 \(q\) 一个整数，因此 \(x^n = (x^p)^q = e\) .

\section{群在集合上的作用}

\subsection{幺半群作用于集合}

定义1. 设M为一个幺半群，其运算用乘法表示，单位元记为e，E为一个集合。一个作用 \(\alpha \mapsto f_{\alpha}\) M在E上的作用称为M在E上的左（相应地，右）作用，如果 \(f_{e} = \mathrm{Id}_{\mathrm{E}}\) 和 \(f_{\alpha \beta} = f_{\alpha} \circ f_{\beta}\) （相应地 \(f_{\alpha \beta} = f_{\beta} \circ f_{\alpha}\) ）对所有 \(\alpha, \beta \in \mathbf{M}\) .

换言之，幺半群M在集合E上的左（相应地，右）作用是从M到幺半群 \(E^{E}\) （分别为 \(E^{E}\) 的相反幺半群）与映射的复合。若与M的作用相对应的作用律记为左（分别为右）乘法，则该作用为左（分别为右）运算这一事实可用公式表示为

\[
e. x = x; \alpha . (\beta . x) = (\alpha \beta). x \quad \text { for } \alpha , \beta \in \mathrm{M} \text { and } x \in \mathrm{E}.\tag{1}
\]

\[
(\text { resp. } x. e = x; (x. \alpha). \beta = x. (\alpha \beta) \quad \text { for } \alpha , \beta \in M \text { and } x \in E).
\]

在这些条件下，也称M在E上左（相应地，右）作用，且相应的作用法则为幺半群M在E上的左（相应地，右）运算律。

设 M 是一个幺半群；带有 M 在 E 上的左（相应地，右）作用的集合 E 称为左（相应地，右）M-集。若映射 \(\alpha \mapsto f_{\alpha}\) 的M到 \(E^{E}\) 是单射。

设 E 为一个集合；其典范作用为 \(E^{E}\) 在E（§3，第1号，例3）上是左作用。

\begin{enumerate}
\def\labelenumi{(\arabic{enumi})}
\setcounter{enumi}{1}
\tightlist
\item
  设 M 是一个幺半群。由 M 上的运算导出的 M 对自身的左（相应地，右）作用（§ 3，第 3 号，例 7）是 M 对自身的一个左（相应地，右）运算。在考虑这一运算时，我们说 M 通过左（相应地，右）平移作用于自身。
\end{enumerate}

设 E 是一个左（相应地，右）M-集合， \(M^{0}\) M 的对偶幺半群。在相同作用下，幺半群 \(M^{0}\) 在 E 上右（相应地，左）作用。所得的 \(M^{0}\) -集合称为 M-集合 E 的对偶集合。关于左 M-集合的定义和结论可平移至右 \(M^{0}\) -集，当过渡到相反结构时。

在本段其余部分中，除非另有说明，我们将仅考虑左M-集，并简称为M-集。其作用法则将用左乘法表示。

设E为一个集合。设G为作用于E上的一个群。对于所有 \(\alpha\) 在 G 中，元素 \(E^{E}\) 上的作用，定义为 \(\alpha\) 是 E 的一个置换（ \(\S\) 2，第 3 号，例 2）。给定 G 在 E 上的一个作用，因此等价于给定一个从 G 到 \(S_{E}\) .

依照第3节第3项，我们作出如下定义：

定义2. 设 M 为一个幺半群，且 E 与 E' 为 M-集合。称映射 f 从 E 到 E' 满足，对所有 \(x \in E\) 以及所有 \(\alpha \in M\) , \(f(\alpha \cdot x) = \alpha \cdot f(x)\) 被称为M-集同态（或M-态射，或与M的运算相容的映射）。

M-集合的恒等映射是M-态射。两个M-态射的复合仍是M-态射。一个M-集合到另一个M-集合的映射成为同构的充要条件是：它是一个双射的M-态射，并且其逆映射也是M-态射。

设 \((\mathbf{E}_i)_{i\in \mathbf{I}}\) 设为一族M-集合，并令 \(\mathbf{E}\) 为 \(\mathbf{E}_i\) 。幺半群 M 作用于 \(\mathbf{E}\) 中条件的扩张 \(\alpha .(x_i)_{i\in \mathbf{I}} = (\alpha .x_i)_{i\in \mathbf{I}}\) 和 \(\mathbf{E}\) ，在此作用下，是一个M-集；设 \(\mathbf{E}'\) 是一个M-集；映射 \(f\) 于 \(\mathbf{E}'\) 到 \(\mathbf{E}\) 是M-态射当且仅当 \(\mathbf{pr}_i\circ f\) 是 \(\mathbf{E}'\) 到 \(\mathbf{E}_i\) 为所有 \(i\in \mathbf{I}\) .

的M-态射。设E是一个M-集，F是E在M作用下稳定的子集；在诱导的运算下，F是一个M-集，且典范嵌入 \(F \rightarrow E\) 是M-态射。设E是一个M-集，R是E上与M作用相容的等价关系；商集E/R在商作用下是一个M-集，且典范映射

是幺半群同态， \(E \rightarrow E/R\) 是M-态射。设E是一个M-集，R是E上与M作用相容的等价关系；商集E/R在商作用下是一个M-集，且典范映射

设 \(\phi: \mathbf{M} \to \mathbf{M}'\) -集和 \(\mathbf{E}\) 一个 \(\mathbf{M}\) -集。映射 \(\mathbf{E}'\) 一个 \(\mathbf{M}'\) 满足对所有 \(f\) 于 \(\mathbf{E}\) 到 \(\mathbf{E}'\) 有 \(x \in \mathbf{E}\) 和 \(\alpha \in \mathbf{M}\) ,

\[
f (\alpha . x) = \phi (\alpha). f (x)
\]

称为 \(\phi\) -从E到E'的态射（参见§3，第1号）。

运算律的推广。给定（例如）三个集合 \(F_{1}, F_{2}, F_{3}\) ，排列 \(f_{1}, f_{2}, f_{3}\) 于 \(F_{1}, F_{2}, F_{3}\) 分别，以及基集上的一个阶梯F \(F_{1}, F_{2}, F_{3}\) （集合论，IV，§ 1，第 1 号），我们可以沿着分层结构 F 的构造逐步定义 F 的一个置换，称为 \(f_{1}, f_{2}, f_{3}\) 至F（集合论，IV，§1，第2号）；我们将其记为 \(\phi_{F}(f_{1}, f_{2}, f_{3})\) .

那么设 G 为一个群，且 \(h_{i}\) G 到对称群的一个同态， \(F_{i}\) (i = 1, 2, 3)，换言之，G 在 \(F_{i}\) 该映射 \(x \mapsto x_{F} = \phi_{F}(h_{1}(x), h_{2}(x), h_{3}(x))\) 是G到 \(S_{F}\) ，换言之，即G在F上的一个作用，称为 \(h_{1}, h_{2}, h_{3}\) 到F。设P是F的一个子集，使得对于所有 \(x \in G\) , \(x_{F}(P) = P\) ；令 \(x_{P}\) 为 \(x_{F}\) 到

P；则映射 \(x \mapsto x_{P}\) 是G在P上的一个运算，也称为 \(h_{1}, h_{2}, h_{3}\) 到P。

例如，设 \(\mathbf{K}\) 和 \(\mathbf{L}\) 是上的两个层；取 \(\mathbf{F}_1, \mathbf{F}_2, \mathbf{F}_3\) 为 \(\mathbf{F}\) 的子集构成的集合，取 \(\mathbf{K} \times \mathbf{L}\) 和 \(\mathbf{P}\) 为 \(\mathbf{K}\) 到 \(\mathbf{L}\) 的映射构成的集合，与它们的图等同。如果 \(w \in \mathbf{P}\) 和 \(x \in \mathbf{G}\) , \(x_{\mathbf{P}}(w)\) 是映射 \(k \mapsto x_{\mathbf{L}}(w(x_{\mathbf{K}}^{-1}(k)))\) 于 \(\mathbf{K}\) 到 \(\mathbf{L}\) .

\subsection{稳定化子，固定子}

定义3. 设M是作用在集合E上的幺半群，且A和B是E的子集。集合 \(\alpha \in \mathbf{M}\) 的集合，使得 \(\alpha \mathbf{A} \subset \mathbf{B}\) （相应地 \(\alpha \mathbf{A} = \mathbf{B}\) ）称为从A到B的迁移子（相应地，严格迁移子）。从A到A的迁移子（相应地，严格迁移子）称为A的稳定化子（相应地，严格稳定化子）。集合 \(\alpha \in \mathbf{M}\) 的集合，使得 \(\alpha a = a\) 为所有 \(a \in \mathbf{A}\) 称为A的固定子。

一个元素 \(\alpha\) 称M稳定（分别地，严格稳定，分别地，固定）E的子集A，如果 \(\alpha\) 属于 A 的稳定化子（分别地，严格稳定化子，分别地，固定子）。M 的子集 P 被称为稳定（分别地，严格稳定，分别地，固定）E 的子集 A，如果 P 的所有元素都稳定（分别地，严格稳定，分别地，固定）A。A 的固定子包含在 A 的严格稳定化子中，而后者本身又包含在 A 的稳定化子中。

命题1. 设M是作用在集合E上的一个幺半群，且A是E的一个子集。

\begin{enumerate}
\def\labelenumi{(\alph{enumi})}
\item
  A的稳定子、严格稳定子和固定子都是M的子幺半群。
\item
  设 \(\alpha\) 是 \(\mathbf{M}\) ；若 \(\alpha\) 属于严格稳定化子（相应地，固定化子） \(\mathbf{A}\) ，则 \(\alpha^{-1}\) .
\end{enumerate}

设 \(e\) 是 \(\mathbf{M}\) ；于是 \(ea = a\) 对于每个元素 \(a \in \mathbf{A}\) ，从而 \(e\) 属于固定化子 \(\mathbf{A}\) 的一个映射。设 \(\alpha\) 和 \(\beta\) 是...的元素 \(\mathbf{E}\) 稳定 \(\mathbf{A}\) 。则 \((\alpha\beta)\mathbf{A} = \alpha(\beta\mathbf{A}) \subset \alpha\mathbf{A} \subset \mathbf{A}\) 因此， \(\mathbf{A}\) 中的像一起构成 \(\mathbf{M}\) 的稳定化子。类似地，对于 \(\mathbf{A}\) ，由此得 (a)。若 \(\alpha\mathbf{A} = \mathbf{A}\) ，则 \(\mathbf{A} = \alpha^{-1}(\alpha\mathbf{A}) = \alpha^{-1}\mathbf{A}\) 如果对于所有 \(a \in \mathbf{A}\) , \(\alpha a = a\) ，则 \(a = \alpha^{-1}(\alpha a) = \alpha^{-1}a\) ，由此得(b)。

推论。设G是作用在集合E上的群，A是E的子集。A的严格稳定子S和固定子F是G的子群，且F是S的正规子群。

第一个断言由命题得出，且F是S到 \(S_{A}\) 的同态映射的核，该同态映射与S在A上的作用相关联。

群G忠实作用于集合E，当且仅当E的固定子由G的单位元组成。E的固定子是G到 \(S_{E}\) 的给定同态映射的核；该同态映射是单射，当且仅当其核由单位元组成（§4，第5号，定理3）。

设 M 为一个幺半群，E 为一个 M-集合，a 为 E 中的一个元素。a 的固定子、严格稳定子和稳定子分别为 \(\{a\}\) 相等；这个幺半群同样被称为 a 的固定子或稳定子。E 的子集 A 的固定子是 A 中各元素固定子的交集。若 a 的固定子为幺半群 M，则称 a 为 E 的不变元素。若 E 中每个元素都是不变的，则称 M 在 E 上平凡作用。

设 G 是一个作用在集合 E 上的群，并且对于所有 \(x \in E\) ，设 \(S_x\) 为 \(x\) 的分拆集合的双射。对于所有 \(\alpha \in G\) , \(S_{\alpha x} = \alpha S_x \alpha^{-1}\) .

上定义了一个处处有定义的合成律。若 \(s \in S_x\) ，则 \(\alpha s \alpha^{-1}(\alpha x) = \alpha s x = \alpha x\) 的自同态，因此 \(\alpha S_x \alpha^{-1} \subset S_{\alpha x}\) 。如 \(x = \alpha^{-1}(\alpha x)\) , \(\alpha^{-1}S_{\alpha x} \alpha \subset S_x\) 的自同态，因此 \(S_{\alpha x} \subset \alpha S_x \alpha^{-1}\) .

类似地可以看出，若A和B是E的两个子集，且T是从A到B的转移子（相应地，严格转移子），则从 \(\alpha A\) 到 \(\alpha B\) 等于 \(\alpha T\alpha^{-1}\) .
